\maketitle
\noteprovenance

\begin{abstract}
\noindent
Erd\H{o}s Problem~\#1049 asks whether
\(F(t)=\sum_{n\ge1}(t^n-1)^{-1}\) is irrational for every rational \(t>1\).
We prove that \(F(a/b)\) is irrational for every pair of coprime integers
\(a>b\ge1\) with \(b^{\mu}<a\), where \(\mu=C_1/C_0=2.464978683574975\ldots\),
equivalently \(\log b/\log a<\theta^{*}=1/\mu=0.40568302138406054\ldots\), and
where \(C_1=1091/2\) and \(C_0=266-(3/\pi^{2})(225-J)\) are the constants
printed by Zudilin in 2004 for his \((14,12,14;27)\) direction, whose ratio is
the integer-base irrationality-exponent bound proved there.  The first base this settles beyond the
published region of Bundschuh and V\"a\"an\"anen is \(31/4\), together with
every power \((31/4)^{r}\), \(r\ge1\).

The region theorem is an ordinary proof.  It cites the polynomial conclusion of
Zudilin's Lemma~7 together with the inputs of that lemma's own proof, and every
remaining step is proved here: the exact degrees, the cyclotomic limit,
positivity, the homogenisation and the denominator balance.  The finite
arithmetic is checked by the Lean kernel: the rational bracket
\(81/200<\theta^{*}<1/2\) around the derived constant, the integer identities of
the degree computation, and the comparisons that place \(31/4\) inside the
region and outside the earlier one.

A degree-budget theorem bounds the reach of the same mechanism.  For a family
of integer polynomial pairs with decay exponent \(\sigma\), degree exponent
\(\delta\) and coefficient-height exponent \(h\), all three independent of the
base, one has \(\sigma\le\delta\), so the homogenised forms are proved to decay
only below \(\sigma/(\sigma+\delta)\le1/2\).  The base \(3/2\) sits at
\(\log2/\log3=0.6309297535714574\ldots\), which is \(0.2252467\ldots\) above
\(\theta^{*}\).

We also determine the \(q\)-order of Zudilin's normalised Hankel determinant
exactly: \(\operatorname{ord}_qV_N^{*}=N(N-1)(2N-1)/6\) at every rank, with
leading coefficient \((N!)^{2}(N+1)!/2^{N}\), where the source proves the
inequality alone.  The initial monomial of every transformed row is identified
here by a filtered reciprocal-state argument.  The all-depth recurrence, the
initial monomial of the first transformed row and the closed-form assembly are
kernel-checked; the Lean formalisation of the row identity stops at the second
row, so the Lean identification of the determinant itself retains a row
hypothesis.

At \(3/2\) we prove exact obstructions to several approximation methods and a
conditional reduction of the number of rows needed to cancel endpoint residues.
The integer bracket \(2^{64}<3^{41}<2^{65}\) gives a uniform gap greater than
\(3/13\) between \(\log2/\log3\) and every threshold in the rectangular
Hermite--Pad\'e exponent model studied here.  The scalar factor and its
universally forced first-order border also fall short of the required
\(39/41\) charge.  These comparisons exclude the stated height mechanisms.

For the additive alternative, the same bracket gives a four-jet collision
among binary selectors whenever \(M\ge130T+2S\) polynomial pairs are used at
bottom depth \(41T\) and top depth \(S\).  The counting threshold is exact
at \(T=1\).  If the rows modulo \(2^S3^R\) are unimodular and their adjacent
determinants vanish, a B\'ezout shear reduces all selector sums to one residue
coordinate; \(S+2R\) rows then suffice for \(R>0\).  Neither coordinate need
be invertible.  A separate bounded-fibre theorem states the additional
multiplicity estimate that would force a collision with distinct analytic
remainders.

Integer scalar content cannot improve the ratio of a determinant's local
divisor to its absolute height, while coordinatewise clearing at \(3/2\)
fails an exact exponential growth bound.  The unresolved step at that base is
an actual approximation family combining primitive coefficients, a nonzero
remainder, and decay faster than height growth.  No irrationality theorem at
\(3/2\) is proved.
\end{abstract}

\begin{paperresult}{What the paper proves}
\textbf{New irrational bases.} $F(a/b)$ is irrational for coprime
$a>b\ge1$ with $\log b/\log a<\theta^{*}=0.40568302138406054\ldots$
(Theorem~\ref{long1049:res:region}).  The base $31/4$ and every power $(31/4)^{r}$ lie
in that region and outside the published region of Bundschuh and
V\"a\"an\"anen, whose cutoff is $1/2-1/\pi^{2}=0.398678816\ldots$
(Theorem~\ref{long1049:res:31over4}).  \textbf{Evidence.} An ordinary proof citing
Zudilin's Lemma~7 and the inputs of its proof, with the finite height
arithmetic at $31/4$ checked by the Lean kernel.  \textbf{Reach of the
mechanism.} A degree-budget theorem caps every base-uniform family of this
shape at $\sigma/(\sigma+\delta)\le1/2$, and $3/2$ sits at
$0.6309297535714574\ldots$ (Theorem~\ref{long1049:res:archcap}).
\textbf{Exact arithmetic at $3/2$.} The bracket $2^{64}<3^{41}<2^{65}$ excludes
the stated scalar and first-order Hermite--Pad\'e height mechanisms there.
It also yields the sharp finite four-jet count at bottom depth $41$, with a
B\'ezout--Pl\"ucker reduction under explicit minor and unimodularity hypotheses.
\textbf{Transfer obstruction.} Direct integer-base clearing fails by an exact
growth inequality and forcing recurrence.  \textbf{Open boundary.} No
primitive noncollapsed approximation family with a sufficiently small
nonzero remainder is constructed at $3/2$, so no irrationality at $3/2$ is
proved, and the universal assertion over all rational $t>1$ remains open.
\end{paperresult}

\section{Introduction}\label{long1049:sec:problem}

Let $t>1$ be a rational number and let $\tau(n)$ count the divisors of $n$.
Erd\H{o}s Problem~\#1049 asks whether
\[
 F(t)=\sum_{n\ge1}\frac{1}{t^{n}-1}=\sum_{n\ge1}\frac{\tau(n)}{t^{n}}
\]
is irrational~\cite[p.~102]{erdos1988}.

The two forms agree by expanding
$(t^{n}-1)^{-1}=\sum_{k\ge1}t^{-nk}$ and collecting the terms with the same
exponent, the coefficient of $t^{-m}$ being the number of divisors of $m$.
The question is a conjecture of Chowla; Erd\H{o}s proved it for every integer
$t\ge2$~\cite{erdos1948}.  Bloom's current catalogue record reproduces the
displayed rational-\(t\) question, labels it \emph{open}, attributes it to
Chowla, and points to Erd\H{o}s's 1988 statement on p.~102 and the 1948
integer-base theorem~\cite{erdosproblems}.  This note proves $F(a/b)$
irrational for every pair of coprime integers $a>b\ge1$ with
$\log b/\log a<\theta^{*}=0.40568302138406054\ldots$, which settles the base
$31/4$ and every power $(31/4)^{r}$, $r\ge1$.  The universal conjecture over
all rational $t>1$ remains open; individual non-integral rational bases are
known, including $7/2$ below.

Write $t=r/s$ in lowest terms with $r>s\ge1$, so that $s=1$ is exactly the
integer case Erd\H{o}s settled.  The resistant explicit base of least naive
height $H(r/s)=\max(r,s)$ is $t=3/2$.  A published height criterion of Bundschuh and
V\"a\"an\"anen~\cite[Thm.~2, p.~177; hypotheses pp.~175--176]{bv1994}
settles a family of rational bases restricted by a
height condition; that family contains $7/2$ and does not contain $3/2$.
Between the two lies the question this note is about: what exactly stops the
integer-base argument from running at $3/2$?

\paragraph{Relation to prior work.}
The \emph{Formal Conjectures} file for Problem~\#1049 contains the conjecture
and integer-base theorem as \texttt{sorry} placeholders, but it also proves
the Lambert identity between the two displayed series for rational
$t$~\cite{formalconjectures1049}.  In the present $t>1$ regime its
\texttt{lambert\_convergent} branch is an ordinary convergent-series proof;
the same file's $|t|\le1$ branch instead uses Lean's convention that the
\texttt{tsum} of a nonsummable series is zero.  Thus it supplies genuine
checked prior art for the identity, but no irrationality theorem.  The present
note formalises propositions about the clearing argument and likewise does
not answer the conjecture.

Erd\H{o}s's positive-integer theorem of
1948~\cite{erdos1948} sits inside a larger integer-base literature.
At the level of functions, Rivin proves that if a sequence \(\gamma\) and its
divisor-sum sequence are both eventually linearly recurrent, then \(\gamma\)
is finitely supported
\cite[Theorem~1.1, p.~2; proof pp.~6--7]{rivin2026}.
His periodic-coefficient corollary therefore shows that
\[
  \sum_{n\ge1}\frac{z^n}{1-z^n}
\]
is not a rational function~\cite[Corollary~6.4, p.~9]{rivin2026}.
This is an exact structural statement about the function underlying
Problem~\#1049, but it gives no irrationality statement for a special value
at \(z=1/t\): a nonrational function may take rational values at particular
rational points.

In 1991
Borwein proved the irrationality of shifted series
$\sum_{n\ge1}(t^{n}+w)^{-1}$ at integer bases $t\ge2$, for every nonzero
rational $w$ with $w\ne-t^m$ for all $m\ge1$, by Pad\'e approximation rather
than by digit clearing; his estimates also show that these values are not
Liouville numbers~\cite[Thm.~4, pp.~257--258]{borwein1991}.  In 2001 Van
Assche recovered the integer-base irrationality and the bound
$\mu(F(p))\le 2\pi^2/(\pi^2-2)=2.50828\ldots$ using little $q$-Legendre
Pad\'e approximants
~\cite[Thm.~1, p.~10; proof pp.~10--11]{vanassche2001}.  His more general
Theorem~3 proves irrationality of
$\sum_{k\ge1}(cp^k-1)^{-1}$ for an integer $p>1$ and fixed rational $c$
away from the poles~\cite[Thm.~3, p.~14]{vanassche2001}; it does not cover a
rational noninteger base $t$ in $F(t)$, because the multiplier needed to
write $t^k$ over an integer base varies with $k$.

The same Lambert value was already the target of Amdeberhan and
Zeilberger's $q$-WZ construction~\cite{az1998}.  Here $p>1$ is the integer-base
parameter and the little-$q$ kernel uses $q=p^{-1}$.  The two constructions
share that bivariate little-$q$-Legendre Pad\'e kernel, but Van Assche's
diagonal does not satisfy the Amdeberhan--Zeilberger scalar recurrence.
Their Theorems~1 and~2 do prove, respectively, the irrationality of the
non-alternating $h_p(1)$ and alternating $q$-logarithm values for their stated
integer parameters, with reported irrationality measure $4.80$; those are
external integer-parameter results and do not transfer to the rational
noninteger base $3/2$ studied here.

Amdeberhan--Zeilberger use the moving diagonal
$P_n(p^{n+1}\mid p^{-1})$, whereas Van Assche uses
$P_n(p^n\mid p^{-1})$.  Van Assche also records, citing
Borwein's 1992 Lemma~2, the neighbouring evaluation
$P_{n-1}(c p^{n+1}\mid p^{-1})$~\cite{borwein1992,vanassche2001}.
The shared kernel therefore does not license transfer of recurrence,
endpoint, lattice, or valuation claims between the diagonals.  Indeed, if
$A_n(p)=P_n(p^n\mid p^{-1})$, exact substitution at $n=0$ into the
Amdeberhan--Zeilberger operator leaves
\[
 -p(p-1)^2(p+1)(p^5+2p^4+2p^3+2p^2+2),
\]
which is nonzero for every real $p>1$.  One nonzero residual is decisive for
non-transfer of that recurrence.

Lean checks the
\lword{Erdos1049/QAperyDiagonalNonEquivalence.lean}{67}{vanAsscheQAperyResidualAtZero_eq}{exact
residual factorisation} and its
\lword{Erdos1049/QAperyDiagonalNonEquivalence.lean}{94}{vanAsscheQAperyResidualAtZero_ne_zero}{nonvanishing
for $p>1$}.  These are finite statements at $n=0$; they do not supply a
recurrence for either diagonal.  No recurrence, endpoint, lattice, or
valuation statement for one diagonal is used for the other; the displayed
exact residual is the sole claim made here about their incompatibility.

There is a second distinction at rational base $3/2$.  Over the checked range,
the raw rational approximation errors decrease while the selected cleared
integer forms grow from $n\ge2$.  This says that the chosen integerisation does
not produce small linear forms; it does \emph{not} say that the raw approximants
fail to converge.

In 2013 Vandehey proved that
$\sum_{n\ge1}d(n)a_n/b^n$ is irrational whenever $b>1$ is an integer and
$(a_n)$ ranges in a finite integer alphabet excluding zero; taking
$a_n=(-1)^n$ completes the elementary digit method for integer bases
$b\le-2$~\cite[Thm.~1.2, p.~2]{vandehey2013}.  His companion theorem permits
a finite alphabet of nonnegative integers containing zero, provided the
coefficient sequence is not eventually zero
\cite[Thm.~1.1, p.~2]{vandehey2013}.  All
three retain an integer base.  Luca and Tachiya's periodic-coefficient theorem
likewise proves irrationality at every integer base of absolute value greater
than one; their full-support example strengthens this to joint linear
independence for finite families of iterated divisor functions
\cite[Theorem~A, p.~139; Example~1, p.~140]{lucatachiya2017}.  This is strong
integer-base evidence, but it does not cover a rational noninteger base.

In 2004 Zudilin obtained the uniform bound
$\mu(F(p))\le2.46497868\ldots$ for every integer
$p\notin\{0,\pm1\}$, by way of Heine's basic transform and a permutation
group~\cite[Thm.~1, p.~154; Secs.~4--5, pp.~159--162]{zudilin2004}.
The ordinary-hypergeometric antecedent is Rhin and Viola's $S_5$ action and
twelve-coset denominator reduction for rational forms in
$\zeta(2)$~\cite[Sec.~3, pp.~38--42; Sec.~4, pp.~46--51]{rhinviola1996}.

Those sources motivate the permutation, denominator and polynomial-
specialisation architecture of Section~\ref{long1049:sec:endpoints}; the endpoint
lemmas there are abstract and are not yet applied to either source's actual
coefficient family.

The neighbouring problem of Lambert
subseries $\sum_{n\in A}(t^{n}-1)^{-1}$ over a restricted index set $A$ is
treated by Kova\v{c} and Tao~\cite{kovactao2024}.

The rational non-integer
progress relevant here is instead the height criterion of Bundschuh and
V\"a\"an\"anen~\cite{bv1994}, which remains the strongest explicit numerical
threshold among the sources compared here for this value.  Zudilin's later
Pad\'e-and-Hankel treatment also admits non-integral rational bases: its
generalized $q$-logarithm remarks that its results also hold for non-integer
$p=r/s$ with $|p|>1$, under an assumption $\log|r|>c\log|s|$ for some computable
$c>0$~\cite{zudilin2016}.  No value of $c$ is computed there.  That remark
announces the shape of a rational-base extension for the generalized
$q$-logarithm of that paper.  Section~\ref{long1049:sec:region} identifies an admissible
constant for $F$ itself, $c=\mu=2.464978683574975\ldots$, which is the 2004
bound quoted above, by specialising the 2004 forms at a rational base and
carrying out the denominator accounting, and derives the region and its first
new base from that constant.

A third published rational-base region for this same series is Duverney's
Th\'eor\`eme~2, whose hypothesis is
$\log|s|/\log|r|<\frac13\bigl(1-3/\pi^{2}\bigr)=0.2320\ldots$
\cite[Th\'eor\`eme~2, p.~174]{duverney1996}; his Th\'eor\`eme~1 for general
numerators is weaker still.  Both cutoffs lie strictly inside the Bundschuh and
V\"a\"an\"anen region.  Matala-aho, V\"a\"an\"anen and Zudilin's combined
treatment of $q$-logarithms keeps the integer hypothesis $p=1/q\in\Z\mathbin{\backslash}\{0,\pm1\}$
throughout, and states in print that its methods do not sharpen the
$q$-harmonic case of~\cite{zudilin2004}
\cite[abstract p.~879; introduction p.~880]{matalaaho2006}.  So none of the
printed regions for $F$ compared here reaches $\log4/\log31$.

\paragraph{Erd\H{o}s's integer-base argument, in outline.}
For an integer $b\ge2$, Erd\H{o}s rewrites
$F(b)=\sum_{n\ge1}\tau(n)b^{-n}$.  A Chinese-remainder construction forces
arbitrarily long blocks in which the divisor coefficients have the powers of
$b$ needed to make the corresponding base-$b$ digits zero.  Explicit bounds
control the middle and far tails, while positivity proves that the expansion
does not terminate.  The resulting base-$b$ expansion has arbitrarily long
zero blocks without being eventually zero and is therefore irrational
~\cite[pp.~63--66]{erdos1948}.

\paragraph{The direct cut-and-clear attempt studied here.}
A more naive attempt is to suppose $F(b)=p/q$, cut at $N$, and multiply the
remaining identity by $qb^N$.  This does not by itself trap a positive integer
below $1$: the first uncleared tail contribution is
$q\tau(N+1)/b$.  The sections below isolate additional corridor hypotheses
under which a bounded-window version of this attempt would work, and then show
why those hypotheses fail at $3/2$.

Write $\beta=r/s$ for the base and $c(n)$ for the coefficient of $\beta^{-n}$,
so that $c=\tau$ in the case at hand.  The term $c(n)\beta^{-n}$ is
$c(n)s^{n}/r^{n}$.  Clearing the
power of $r$ leaves the numerator factor $s^{n}$ in place.  That factor is
invisible when $s=1$ and grows geometrically when $s\ge2$.  At $3/2$ it is
$2^{n}$.

\paragraph{Terminology.}
The linear-form method uses polynomial pairs before specialisation and integer
pairs afterwards.  We reserve \emph{row} for an integer pair $(U,V)$.  Its
\emph{integer scalar content} is $\gcd(|U|,|V|)$; a row is \emph{primitive}
when this number is $1$, and dividing by it is \emph{primitive normalisation}.
A polynomial pair may instead have a polynomial common factor in $\Z[X]$;
that is a different operation and is not called row content here.  The
\emph{exterior determinant} of two integer rows is
$U_{n}V_{m}-U_{m}V_{n}$, the determinant of the $2\times2$ matrix they form.
Two quantities attached to that determinant are compared throughout: an
integer dividing it, which is a local gain, and its absolute value, which is an
Archimedean cost; we call that comparison the \emph{local-to-Archimedean
balance}.

The \emph{endpoints} of a coefficient polynomial, taken relative to
the declared width $W$ of Section~\ref{long1049:sec:endpoints}, are its constant
coefficient and its coefficient at $W$; we call these the \emph{constant
endpoint} and the \emph{top endpoint}, so the top endpoint is the coefficient
at $W$ and not the leading coefficient unless the two agree.  A \emph{unit}
endpoint is one equal
to $\pm1$, the units of $\Z$; Theorem~\ref{long1049:res:endpoints} is the reason only these
two coefficients decide divisibility by $3$ and by $2$ after specialisation at
$(3,2)$.  A \emph{jet} is a
residue of a specialised coefficient modulo a prime power: the bottom jet is
its residue modulo $3^{R}$ and the top jet its residue modulo $2^{S}$, and $R$
and $S$ are the bottom and top \emph{depths}.  A jet vanishes exactly when the
prime power in question divides the specialised coefficient.

\paragraph{The shortfall at $3/2$.}
The elementary method and the linear-form method are both examined below.  For
the height argument, the exact bracket $2^{64}<3^{41}<2^{65}$ gives a gap greater
than $3/13$ across the whole admissible rectangular exponent cone and fixes the
$39/41$ denominator-charge comparison.  Neither the source scalar factor nor
that factor together with the universally forced first-order residual border
reaches the required charge.  These comparisons leave higher residual
valuations, determinant cancellation, and other integral models untouched.
For the additive argument, the same bracket replaces the generic
$4R+2S$ threshold by $130T+2S$ at depth $R=41T$ and gives an exact first-row
failure at $T=1$.

For
the coordinatewise clearing scheme the leftover at each step is the forcing
term of an exact recurrence, of size at least
$2^{N+1}$ whenever $s\ge2$ and the scaling constant $B$ and the coefficient
$c(N+1)$ are at least $1$
(Theorem~\ref{long1049:res:forcing}), and the
scheme itself is excluded at $3/2$ for every shift $N\ge1$ and every cleared
window of width $K\ge1$ (Theorem~\ref{long1049:res:nocorridor}).

For the linear-form
constructions we examine two possible sources of $2$-adic and $3$-adic gain.
Integer scalar content is exactly neutral, since it scales the exterior
determinant and its absolute height by the same factor
(Theorem~\ref{long1049:res:content}), while unit endpoints keep both $2$ and $3$ out of
any common divisor of the two specialised evaluations
(Theorem~\ref{long1049:res:endpoints} and Proposition~\ref{long1049:res:commonmult}).
Corollary~\ref{long1049:res:nomult} states the two scoped exclusions together; neither
is claimed to be necessary for every linear-form proof.  Separately, the scalar
parameter margin is negative under the assumed source inequality
(Theorem~\ref{long1049:res:scalar}).

One candidate pursued here is additive: an
integer relation among rows that cancels the endpoint jets.
Theorem~\ref{long1049:res:jetkernel} shows that a nonzero relation with coefficients in
$\{-1,0,1\}$ cancelling all four jets exists whenever the bottom depth is
positive and the number of coefficient pairs is at least $4R+2S$.  It does not
show that the resulting combination has a nonzero polynomial pair or a nonzero
remainder.  Problem~\ref{long1049:prob:kernel} gives a precise sufficient specification
for that particular candidate architecture, not a necessary condition for
solving Problem~\#1049.

Sections~\ref{long1049:sec:sevenhalves} and~\ref{long1049:sec:pade} record two external methods
and what each leaves unproved at $3/2$.

\paragraph{Sharpness.}
Two questions of scope are worth isolating.  The corridor bound of
Theorem~\ref{long1049:res:corridorbound} is exponential on the left and linear on the
right, so for a fixed numerator and any $s\ge2$ a corridor can survive only for
bounded $N+K$; the base $3/2$ is the case in which the crossing has already
happened at the smallest admissible window, which is why the exclusion there
holds for all $N\ge1$ and $K\ge1$ with no further restriction.  The height
criterion of~\cite{bv1994} is restricted by a height condition satisfied at
$7/2$ and not at $3/2$, so the two bases are separated by that criterion rather
than by a universal obstruction.  The exponent $65$ in
$3^{41}<2^{65}$ cannot be replaced by $64$.  The separate direct comparison
$2^{129}<3^{82}$ then proves that at bottom depth $41$ the least number of input
rows satisfying the ambient-cardinality inequality is exactly $130+2S$:
$129+2S$ does not satisfy it.  That is a statement about the counting test, and
a family with fewer rows may still have a collision.

For $R=41T$ with $T>1$, the displayed $130T+2S$ bound is uniform and sufficient
and is not always the least such number.

\statusboundary

\paragraph{Results and boundary.}
Theorems~\ref{long1049:res:region} and~\ref{long1049:res:31over4} are ordinary mathematics.  They
consume Zudilin's Lemma~7 together with the inputs of that lemma's own proof,
and every further step is proved in Section~\ref{long1049:sec:region}.  They are not
kernel-checked.  The linked Lean declarations establish exact inequalities,
congruences, finite collision counts, recurrence identities, and exclusions for
the named models; among them are the height comparisons that place \(31/4\)
inside the enlarged logarithmic region and outside the earlier one.  No Lean
declaration in this release's library carries an irrationality statement at a
rational noninteger base.  Bundschuh and V\"a\"an\"anen's irrationality theorem
at \(7/2\), Rivin's functional nonrationality theorem, and the two-variable
Mahler theorem used below remain external results.  At \(3/2\), the missing
step is still an actual primitive approximation family with nonzero remainder
and an asymptotic height margin.

\paragraph{Structure.}
Section~\ref{long1049:sec:region} states and proves the region theorem, the base $31/4$
and the degree-budget cap on the mechanism that produces them.
Section~\ref{long1049:sec:sharp} derives the height and charge exclusions from the sharp
$41/65$ power certificate.  Section~\ref{long1049:sec:primitive} proves that integer scalar content is neutral for the
local-to-Archimedean balance.

Section~\ref{long1049:sec:endpoints} proves the endpoint
congruences at $(3,2)$, deduces from them and from
Section~\ref{long1049:sec:primitive} that neither integer scalar content nor a common divisor
of the two specialised evaluations supplies the targeted endpoint gain, gives the
four-jet collision count and its conditional B\'ezout--Pl\"ucker compression, and records one further
exclusion on Zudilin's scalar parameters.
Sections~\ref{long1049:sec:corridor} and~\ref{long1049:sec:tail} return to the elementary
clearing scheme and record what the residue $s^{n}$ costs there, first as an
exclusion and then as an exact recurrence with a lower bound on the surviving
term.  Sections~\ref{long1049:sec:sevenhalves} and~\ref{long1049:sec:pade} record what is and is
not formalised of two external routes, together with the separate $81/200$
logarithmic comparison.  Section~\ref{long1049:sec:open} separates kernel escape from
asymptotic adequacy and states the remaining obligations.  Linked phrases open the corresponding Lean
declaration at the pinned source revision~\commitshort.

\medskip
\noindent\textbf{Keywords.} irrationality; Lambert series; rational base;
Pad\'e approximation; Lean~4.
\qquad
\textbf{MSC 2020.} 11J72 (primary); 11J82, 68V20 (secondary).

\section{The region \texorpdfstring{$b^{\mu}<a$}{b\textasciicircum mu < a} and the base \texorpdfstring{$31/4$}{31/4}}\label{long1049:sec:region}

This section states the note's leading results.  Fix Zudilin's direction
$(\alpha_0,\alpha_1,\alpha_2;\beta)=(14,12,14;27)$ of
\cite[Sec.~5]{zudilin2004} and put
\[
 C_1=(\alpha_0+\alpha_1+\alpha_2)\beta-\tfrac12(\alpha_1^{2}+\alpha_2^{2}+\beta^{2})
 =\frac{1091}{2},
 \qquad
 C_0=266-\frac{3}{\pi^{2}}\,(225-J),
\]
\[
 J=\sum_{i=1}^{13}\bigl(\psi_1(u_i)-\psi_1(v_i)\bigr),
 \qquad
 \psi_1(x)=\sum_{k\ge0}\frac1{(k+x)^{2}},
\]
the thirteen demi-intervals $[u_i,v_i)$ being
$[1/14,1/12)$, $[1/7,1/6)$, $[3/14,1/4)$, $[2/7,1/3)$, $[5/14,2/5)$,
$[3/7,7/15)$, $[1/2,8/15)$, $[4/7,3/5)$, $[9/14,2/3)$, $[5/7,11/15)$,
$[11/14,4/5)$, $[6/7,13/15)$, $[13/14,14/15)$.  Numerically
\begin{align*}
 J&=77.943184475009095922589567023992\ldots,\\
 C_0&=221.300088165005025124042696692215\ldots,\\
 \mu&:=\frac{C_1}{C_0}=2.464978683574975037454488275535\ldots,\\
 \theta^{*}&:=\frac1\mu=0.405683021384060541015660305576\ldots
\end{align*}
The direction, the thirteen demi-intervals, the values $C_1=545.5$ and
$C_0=221.30008816\ldots$ and the ratio $C_1/C_0=2.46497868\ldots$ are printed at
\cite[p.~162]{zudilin2004}, where that ratio is the irrationality-exponent bound
of \cite[Thm.~1, p.~154]{zudilin2004} for integer bases.  The further digits
displayed above and the reciprocal $\theta^{*}$ are recomputed here from the
same trigamma series.

The proof below needs $J$ in its combinatorial form as well.  Put
\[
 \omega(\xi)=\max\bigl\{0,\;
 \lfloor14\xi\rfloor+\lfloor13\xi\rfloor-\lfloor12\xi\rfloor-\lfloor15\xi\rfloor,\;
 2\lfloor14\xi\rfloor-\lfloor13\xi\rfloor-\lfloor15\xi\rfloor\bigr\},
\]
the weight attached by \cite[(26)]{zudilin2004} to the six-tuple
$c=(13,14,12,14,15,13)$ of the fixed direction.

\begin{lemma}[the weight is an indicator]\label{long1049:res:omega-indicator}
On $[0,1)$ the function $\omega$ takes only the values $0$ and $1$, and
$\omega=1$ exactly on the union of the thirteen demi-intervals listed above.
\end{lemma}

\begin{proof}
Each of the three expressions inside the maximum is constant on every interval
between consecutive elements of
$\{k/c: c\in\{12,13,14,15\},\ 0\le k\le c\}$, since those are the only points at
which one of the four floor functions changes.  There are $48$ such intervals in
$[0,1)$; evaluating $\omega$ at the midpoint of each gives the value $1$ on the
thirteen listed demi-intervals and $0$ elsewhere, and the thirteen listed
intervals are unions of consecutive cells.  The evaluation is exact rational
arithmetic on finitely many points.
\end{proof}

\noindent Writing $\mathbf 1_{[u,v)}$ for indicators, Lemma~\ref{long1049:res:omega-indicator}
and termwise integration of the positive series
$-\psi_1'(\xi)=\sum_{k\ge0}2(k+\xi)^{-3}$ give the two forms of $J$ used below,
\[
 J=\int_0^1\omega(\xi)\sum_{k\ge0}\frac{2}{(k+\xi)^{3}}\,d\xi
  =\sum_{i=1}^{13}\bigl(\psi_1(u_i)-\psi_1(v_i)\bigr),
\]
the first being the integral $\int_0^1\omega\,d(-\psi_1)$ of
\cite[Lemma~2]{zudilin2004}.  Since $\omega$ vanishes near $0$, both are finite.

\begin{theorem}[rational-base region]\label{long1049:res:region}
Let $a>b\ge1$ be coprime integers with
\[
 b^{\mu}<a,
 \qquad\text{equivalently}\qquad
 \frac{\log b}{\log a}<\theta^{*}=0.40568302138406054\ldots
\]
Then $F(a/b)=\sum_{m\ge1}(( a/b)^{m}-1)^{-1}$ is irrational.
\end{theorem}

The proof occupies Sections~\ref{long1049:sec:source-forms}--\ref{long1049:sec:integer-forms}.  It
cites one external statement, the polynomial conclusion of Zudilin's Lemma~7
together with the coefficientwise argument that proves it; everything else is
carried out here.

\subsection{The source forms as polynomial
identities}\label{long1049:sec:source-forms}

Fix $n\ge1$, write $x$ for an indeterminate or a real number greater than $1$
according to context, and put $q=x^{-1}$.  Set
\[
\begin{gathered}
 a_0=14n+1,\quad a_1=12n+1,\quad a_2=14n+1,\\
 \beta_n=27n+2,\quad N=15n,\quad M_n=266n^{2}+34n+1 .
\end{gathered}
\]
The symbol $\beta_n$ is the upper hypergeometric parameter of
\cite[Sec.~2]{zudilin2004}; the letter $b$ continues to denote the denominator
of the base, and $M_n$ is the integer of that paper's~(16).  With
$(z;q)_m=\prod_{j=0}^{m-1}(1-zq^{j})$, the source's summand is
\begin{equation}\label{long1049:eq:positive-source-form}
 H_n(x)=\sum_{t\ge0}
 \frac{(q^{t+1};q)_{a_1-1}}{(q;q)_{a_1-1}}\,
 \frac{(q;q)_{\beta_n-a_2-1}}{(q^{a_2+t};q)_{\beta_n-a_2}}\,q^{a_0t},
\end{equation}
and \cite[(8)--(11)]{zudilin2004} gives $H_n=A_nF-B_n$ with
$A_n,B_n\in\Q(x)$.  The two coefficients can be written down.  Let
$\genfrac{[}{]}{0pt}{}{m}{r}_X$ be the Gaussian binomial polynomial, of degree
$r(m-r)$ and with nonnegative coefficients summing to $\binom{m}{r}$, and for
$a_2\le k\le\beta_n-1$ put
\begin{equation}\label{long1049:eq:c-explicit}
 c_{n,k}(X)=(-1)^{a_1+a_2+k+1}X^{e_{n,k}}
 \genfrac{[}{]}{0pt}{}{k-1}{a_1-1}_X
 \genfrac{[}{]}{0pt}{}{\beta_n-a_2-1}{\beta_n-k-1}_X,
\end{equation}
\[
 e_{n,k}=\frac{a_1(a_1-1)}2-\frac{(\beta_n-a_2)(\beta_n-a_2-1)}2
 +\frac{(\beta_n-k)(\beta_n-k-1)}2 .
\]
Then
\begin{align}
 A_n(X)&=\sum_{k=a_2}^{\beta_n-1}c_{n,k}(X)X^{a_0k},\label{long1049:eq:A-explicit}\\
 B_n(X)&=\sum_{k=a_2}^{\beta_n-1}c_{n,k}(X)X^{a_0k}
 \biggl(\sum_{l=1}^{k-a_1}\frac1{X^{l}-1}
 +\sum_{j=1}^{a_0-1}\frac{X^{-j(k-a_1)}}{X^{j}-1}\biggr).\label{long1049:eq:B-explicit}
\end{align}
Put
\begin{equation}\label{long1049:eq:cyclotomic-products}
 D_N(X)=\prod_{l=1}^{N}\Phi_l(X),\qquad
 \nu_{n,l}=\omega(n/l),\qquad
 \Omega_n(X)=\prod_{l=2}^{N}\Phi_l(X)^{\nu_{n,l}},
\end{equation}
with $\Phi_l$ the $l$th cyclotomic polynomial; by
Lemma~\ref{long1049:res:omega-indicator} every $\nu_{n,l}$ is $0$ or $1$, and
$\nu_{n,l}=0$ for $l>15n$ because the intervals all begin at or above $1/14$.

Zudilin's Lemma~7 \cite[p.~161, (23)]{zudilin2004} applies at the parameters above: the tuple is
$c=n\cdot(13,14,12,14,15,13)$ with maximum $m(c)=15n$ and difference $s(c)=n>0$,
and the source's~(14) holds because $12n+1\le14n+1$ and
$26n+2\le27n+2\le28n+2$.  Its conclusion is the inclusion
\begin{equation}\label{long1049:eq:integer-polynomial-pair}
 U_n:=X^{-M_n}\frac{D_N}{\Omega_n}A_n\in\Z[X],\qquad
 V_n:=X^{-M_n}\frac{D_N}{\Omega_n}B_n\in\Z[X].
\end{equation}
Two points about that citation matter.  First, the conclusion used is the one
in $\Z[X]$, not integrality at integer arguments: an integer-valued polynomial
such as $X(X-1)/2$ shows that the two statements differ.  The source's proof
supplies the polynomial form.  Its Lemma~4 clears the two rational coefficients
separately by Gaussian-binomial identities in $\Z[X]$, its Lemma~5 is Heine's
transform in the parameter $q$, and Lemma~7 removes the cyclotomic factors by
polynomial divisibility, $\Omega_n$ being monic.  All three precede the source's
evaluation at an integer $p$ in its~(24), and all three remain valid with an
indeterminate.  Second, the pair extracted from the source's module inclusion is
unique, because $F$ is not a rational function of $X$: as $h\downarrow0$,
splitting the original Lambert sum $F(e^h)=\sum_{m\ge1}(e^{hm}-1)^{-1}$
at $L=\lfloor1/h\rfloor$ gives an initial sum
$h^{-1}\sum_{m\le L}m^{-1}+O(L)$, since
$y^{-1}-1\le(e^y-1)^{-1}\le y^{-1}$ for $0<y\le1$.
The remaining sum is at most
$e^{-h(L+1)}/((1-e^{-1})(1-e^{-h}))=O(h^{-1})$.  Hence
$F(e^h)=h^{-1}\log(1/h)+O(h^{-1})$, so
$(X-1)F(X)\to\infty$ and $(X-1)^{2}F(X)\to0$ as $X\downarrow1$, which no
rational function does.  Hence the coefficients of $F$ and of $1$ in an identity
over $\Q(X)$ are determined, and~\eqref{long1049:eq:integer-polynomial-pair} identifies
them.  This is a statement about $F$ as a function and says nothing about any
individual value.

\subsection{Exact degrees}\label{long1049:sec:degrees}

Let $d_{n,k}=\deg\bigl(c_{n,k}(X)X^{a_0k}\bigr)$.  The Gaussian-binomial degree
formula gives
\[
 d_{n,k}=e_{n,k}+a_0k+(a_1-1)(k-a_1)+(\beta_n-k-1)(k-a_2)
 =\frac{-k^{2}+80kn+3k-340n^{2}-26n}{2},
\]
\[
 d_{n,k+1}-d_{n,k}=40n+1-k>0\qquad(a_2\le k\le\beta_n-2),
\]
the inequality because $k\le\beta_n-2=27n$ and $27n<40n+1$.  The summand of
top degree is therefore the one at $k=\beta_n-1$ alone, no cancellation occurs
there, and
\begin{equation}\label{long1049:eq:exact-degree}
 \begin{gathered}
 K_n:=\deg A_n=d_{n,\beta_n-1}=\frac{1091n^{2}+81n+2}{2},
 \\
 W_n:=\deg U_n=K_n-M_n+\sum_{l\le N}\varphi(l)-\sum_{2\le l\le N}\nu_{n,l}\varphi(l),
\end{gathered}
\end{equation}
the second equality because $\deg D_N=\sum_{l\le N}\varphi(l)$ and
$\deg\Omega_n=\sum_{2\le l\le N}\nu_{n,l}\varphi(l)$, and because
$U_n$ is a polynomial by~\eqref{long1049:eq:integer-polynomial-pair}, so subtracting
$M_n$ from the degree of $D_NA_n/\Omega_n$ is legitimate.  In particular
$U_n\ne0$.

The three integer identities in this step are checked by the Lean kernel:
$2M_n$ against the source's~(16) is
\lword{Erdos1049/RationalBaseContour.lean}{317}{two_mul_zudilinM}{one
declaration}, the degree step $d_{n,k+1}-d_{n,k}=40n+1-k$ is
\lword{Erdos1049/RationalBaseContour.lean}{337}{twoMulZudilinSummandDegree_succ_sub}{a
second}, its positivity on the whole range $a_2\le k\le\beta_n-2$ is
\lword{Erdos1049/RationalBaseContour.lean}{350}{twoMulZudilinSummandDegree_step_pos}{a
third}, and the top value $2K_n=1091n^{2}+81n+2$ is
\lword{Erdos1049/RationalBaseContour.lean}{344}{twoMulZudilinSummandDegree_top}{a
fourth}.  They decide the arithmetic of this step and no analytic step
anywhere in the note.

For the second coefficient, fix $n$ and let $x\to\infty$.  Every factor
in~\eqref{long1049:eq:positive-source-form} tends to $1$ and the geometric factor
$q^{a_0t}$ makes the sum converge, so $H_n(x)=O(1)$, while
$F(x)=\sum_m\tau(m)x^{-m}=O(x^{-1})$.  From $B_n=A_nF-H_n$ and
$\deg A_n=K_n$ it follows that $B_n(x)=O(x^{K_n-1})$; multiplying by
$x^{-M_n}D_N(x)/\Omega_n(x)$ and using
integrality~\eqref{long1049:eq:integer-polynomial-pair} gives
\begin{equation}\label{long1049:eq:V-degree}
 \deg V_n\le W_n-1 .
\end{equation}
One common homogenising width $W_n$ therefore serves both coefficients.

\subsection{The cyclotomic limit}\label{long1049:sec:cyclotomic-limit}

Write $\Sigma_n=\sum_{2\le l\le N}\nu_{n,l}\varphi(l)$.  The elementary
summatory estimate is $\sum_{l\le y}\varphi(l)=3\pi^{-2}y^{2}+O(y\log(2y))$.
Fix one demi-interval $[u,v)$ of Lemma~\ref{long1049:res:omega-indicator}.  The condition
$\{n/l\}\in[u,v)$ holds exactly on the blocks
\[
 \frac{n}{k+v}<l\le\frac{n}{k+u},\qquad k=0,1,2,\dots,
\]
so the summatory estimate gives, for each fixed $k$,
$n^{-2}\sum_{l\text{ in block }k}\varphi(l)\to3\pi^{-2}\bigl((k+u)^{-2}-(k+v)^{-2}\bigr)$.
The blocks with $k\ge L$ consist of integers $l\le n/(L+u)$, and
$\varphi(l)\le l$, so their total is at most $\tfrac12(n/(L+u))^{2}+O(n)$ and
their normalised contribution is $O(L^{-2})$ uniformly in $n$.  Taking
$n\to\infty$ first and then $L\to\infty$ therefore justifies the infinite block
sum, and summing the thirteen demi-intervals gives
$n^{-2}\Sigma_n\to3J/\pi^{2}$.  With $\sum_{l\le15n}\varphi(l)=3\pi^{-2}225n^{2}+O(n\log n)$,
equation~\eqref{long1049:eq:exact-degree} yields
\begin{equation}\label{long1049:eq:degree-limits}
 \frac{K_n}{n^{2}}\to C_1,
 \qquad
 \frac{K_n-W_n}{n^{2}}=\frac{M_n-\sum_{l\le N}\varphi(l)+\Sigma_n}{n^{2}}
 \to266-\frac{3}{\pi^{2}}(225-J)=C_0,
\end{equation}
and hence $W_n/n^{2}\to C_1-C_0$.  This limit replaces Lemmas~1 and~2 of the
source, which are consequently not part of the citation set.

For the Archimedean side, fix a real $x>1$.  From
$\log\Phi_l(x)-\varphi(l)\log x=\sum_{d\mid l}\mu_{\mathrm{Mob}}(l/d)\log(1-x^{-d})$,
with $\mu_{\mathrm{Mob}}$ the M\"obius function, each summand has absolute value
at most $-\log(1-x^{-1})$, and $\sum_{l\le N}\tau(l)=O(N\log(2N))$.  Since
$\nu_{n,l}\in\{0,1\}$,
\begin{equation}\label{long1049:eq:cyclotomic-size}
 \log\frac{D_N(x)}{\Omega_n(x)}
 =\Bigl(\sum_{l\le N}\varphi(l)-\Sigma_n\Bigr)\log x+O_x\bigl(N\log(2N)\bigr),
\end{equation}
whose error term is $O_x(n\log n)=o(n^{2})$.  By~\eqref{long1049:eq:exact-degree} the
bracket equals $W_n-K_n+M_n$.

\subsection{Positive integral forms and the
contradiction}\label{long1049:sec:integer-forms}

Fix the base $x=a/b>1$ and write $P_q=(q;q)_\infty>0$.  Each of the four finite
$q$-Pochhammer products in~\eqref{long1049:eq:positive-source-form} is a product of
factors $1-q^{j}$ with $j\ge1$, hence lies in $[P_q,1]$, so each of the two
ratios lies in $[P_q,P_q^{-1}]$, every summand is positive, and
\[
 P_q^{2}\le H_n(x)\le\frac{P_q^{-2}}{1-q^{a_0}}\le\frac{P_q^{-2}}{1-q},
\]
whence $H_n(x)>0$ and $\log H_n(x)=O_x(1)$ uniformly in $n$.  Cyclotomic
polynomials are positive on $(1,\infty)$, so
by~\eqref{long1049:eq:integer-polynomial-pair}
\[
 \Lambda_n(x):=U_n(x)F(x)-V_n(x)=x^{-M_n}\frac{D_N(x)}{\Omega_n(x)}H_n(x)>0 .
\]
By~\eqref{long1049:eq:exact-degree} and~\eqref{long1049:eq:V-degree} both polynomial degrees are
at most $W_n$, so
\[
 \widehat U_n=b^{W_n}U_n(a/b),\qquad \widehat V_n=b^{W_n}V_n(a/b)
\]
are integers, and $\widehat\Lambda_n:=b^{W_n}\Lambda_n(a/b)
=\widehat U_nF(a/b)-\widehat V_n$ is positive and lies in $\Z F(a/b)+\Z$.  This
is the only non-Archimedean step, and it is where the denominator is paid for.
Combining the size estimates,
\[
 \begin{aligned}
 \log\widehat\Lambda_n
 &=W_n\log b-M_n\log x+\Bigl(\sum_{l\le N}\varphi(l)-\Sigma_n\Bigr)\log x+O_x(n\log n)\\
 &=K_n\log b-(K_n-W_n)\log a+o(n^{2}),
\end{aligned}
\]
the second equality by $\log x=\log a-\log b$ and~\eqref{long1049:eq:exact-degree}.
With~\eqref{long1049:eq:degree-limits} this proves
\begin{equation}\label{long1049:eq:final-limit}
 \lim_{n\to\infty}\frac{\log\widehat\Lambda_n}{n^{2}}=C_1\log b-C_0\log a .
\end{equation}

\begin{proof}[Proof of Theorem~\ref{long1049:res:region}]
The hypothesis $\log b/\log a<\theta^{*}=C_0/C_1$ makes the right side
of~\eqref{long1049:eq:final-limit} negative, so $\widehat\Lambda_n>0$ and
$\widehat\Lambda_n\to0$.  Suppose $F(a/b)=P/Q$ with integers $P$ and $Q\ge1$.
Then $Q\widehat\Lambda_n=P\widehat U_n-Q\widehat V_n$ is a positive integer for
every $n$ and tends to $0$, which is impossible.
\end{proof}

\noindent Above the threshold~\eqref{long1049:eq:final-limit} says that these same
homogenised forms grow like $e^{cn^{2}}$ with $c>0$, and at equality the limit
decides nothing.  Both statements are about the displayed family.

\noindent The hypothesis $b\ge1$ admits $b=1$, where the statement reduces to
the known integer-base theorem.  Coprimality is a restriction on the pairs
treated and is used nowhere in the argument.  Negative bases are excluded,
because Step~4 is a positivity argument.

\begin{theorem}[the first base beyond the published
region]\label{long1049:res:31over4}
$F(31/4)$ is irrational, and so is $F\bigl((31/4)^{r}\bigr)$ for every integer
$r\ge1$.  Here
\[
 \frac{\log4}{\log31}=0.4036981731641997\ldots<\frac{81}{200}<\theta^{*},
\]
while $4^{\mu}=30.483515\ldots<31<4^{\mu_{\mathrm{BV}}}=32.369642\ldots$ with
$\mu_{\mathrm{BV}}=2\pi^{2}/(\pi^{2}-2)=2.508284761994\ldots$, so $31/4$ lies
outside the region $\log b/\log a<1/2-1/\pi^{2}=0.3986788163576622\ldots$ of
\cite[Thm.~2, p.~177]{bv1994} and inside the region of
Theorem~\ref{long1049:res:region}.
\end{theorem}

\begin{proof}
The comparison $\log4/\log31<81/200$ is the integer certificate
$4^{200}<31^{81}$, checked as
\lword{Erdos1049/ZudilinHeightRegion.lean}{26}{thirtyoneFour_power_certificate}{the
power certificate}, and the membership it yields is
\lword{Erdos1049/ZudilinHeightRegion.lean}{45}{thirtyoneFour_mem_zudilinHeightRegion}{checked};
the ratio $\log b/\log a$ is invariant under $(a,b)\mapsto(a^{r},b^{r})$, which
gives the
\lword{Erdos1049/ZudilinHeightRegion.lean}{59}{thirtyoneFour_power_mem_zudilinHeightRegion}{power
family}, and $(31^{r},4^{r})$ are coprime.  The exclusion from the earlier
region is
\lword{Erdos1049/ZudilinHeightRegion.lean}{91}{thirtyoneFour_outside_bundschuhVaananenHeightRegion}{checked}.
That exclusion also has a two-line rational certificate: $31^{2}<4^{5}$ gives
$\log4/\log31>2/5$, and $\pi^{2}<10$ gives $1/2-1/\pi^{2}<2/5$, so
\[
 \frac12-\frac1{\pi^{2}}<\frac25<\frac{\log4}{\log31}<\frac{81}{200}<\theta^{*}.
\]
The remaining comparison $81/200<\theta^{*}$ is proved in
Section~\ref{long1049:sec:regionbracket}.  Theorem~\ref{long1049:res:region} then applies.
\end{proof}

\subsection{The rational bracket around \texorpdfstring{$\theta^{*}$}{theta*}}\label{long1049:sec:regionbracket}

The constant $\theta^{*}$ is defined by a trigamma series, so a decidable
membership test needs explicit rational bounds on it.  Both are finite rational
arithmetic, and neither uses a decimal expansion.

For the lower bound, keep only the $k=0$ term of each of the thirteen
differences $\psi_1(u_i)-\psi_1(v_i)$.  All later terms are positive, so
\[
 J\ \ge\ \sum_{i=1}^{13}\Bigl(\frac1{u_i^{2}}-\frac1{v_i^{2}}\Bigr)
  =\frac{2015640690251}{25971865920}
  >\frac{776}{10}.
\]
Since $\pi>157/50$ and $225-J<225$,
\[
 C_0>266-\frac{3\,(225-776/10)}{(157/50)^{2}}=\frac{5451134}{24649}
 >\frac{88371}{400}=\frac{81}{200}\,C_1,
\]
so $81/200<\theta^{*}$.  For the upper bound the thirteen demi-intervals are
disjoint and ordered and $\psi_1$ is positive and decreasing, so the sum
telescopes below its first term:
\[
 J<\psi_1(1/14)=196+\sum_{k\ge1}\frac1{(k+1/14)^{2}}<196+\frac{\pi^{2}}6<198<225,
\]
whence $C_0<266$ and $\theta^{*}<532/1091<1/2$.  So
\[
 \frac{81}{200}<\theta^{*}<\frac12 .
\]

Both bounds, the definitions of $C_0$, $C_1$ and $\theta^{*}$ from the trigamma
series, the membership of $31/4$ and of every power of it, and the exclusion of
$3/2$ are checked by the Lean kernel in the module
\texttt{RationalBaseContour}, which the library root imports at the pinned
source revision~\commitshort.  The lower bound on $J$ is
\lword{Erdos1049/RationalBaseContour.lean}{159}{zudilinJ_ge}{one declaration},
the rational bound on $C_0$ is
\lword{Erdos1049/RationalBaseContour.lean}{180}{zudilinC0_gt}{a second}, the
comparison $81/200<\theta^{*}$ is
\lword{Erdos1049/RationalBaseContour.lean}{201}{eightyOne_twoHundredths_lt_zudilinContour}{a
third}, the bound $\theta^{*}<1/2$ is
\lword{Erdos1049/RationalBaseContour.lean}{251}{zudilinContour_lt_half}{a
fourth}, the membership of $31/4$ and of its powers are
\lword{Erdos1049/RationalBaseContour.lean}{266}{thirtyoneFour_mem_zudilinContourRegion}{a
fifth} and
\lword{Erdos1049/RationalBaseContour.lean}{279}{thirtyoneFour_power_mem_zudilinContourRegion}{a
sixth}, and the exclusion of $3/2$ is
\lword{Erdos1049/RationalBaseContour.lean}{296}{threeHalves_outside_zudilinContourRegion}{a
seventh}.  Each is a finite comparison against the defined constant; none
of them is any part of the analytic argument of
Theorem~\ref{long1049:res:region}.  A finite-precision evaluation of the trigamma series
with the integral tail bound over its first thousand terms gives the sharper
enclosure $0.4056830211<\theta^{*}<0.4056830214$; that enclosure is a numerical
receipt and is not kernel-checked.

\subsection{Evidence, attribution and the remaining obligation}

\paragraph{Evidence decomposition.}
Theorems~\ref{long1049:res:region} and~\ref{long1049:res:31over4} are ordinary proofs.  Their
citation set is Zudilin's Lemma~7 in its polynomial reading, together with the
inputs of that lemma's own proof, namely his Lemma~3, the exponent $M(a;b)$
of~(16) proved by his Lemma~4, Heine's transform as his Lemma~5, and the
identity~(9)--(11) \cite{zudilin2004}.  Everything else above is proved here:
the nonrationality of $F$ as a function, the transfer of Lemma~7 to explicit
integer polynomials, the exact degree computation, the degree bound on $V_n$,
positivity, the Archimedean size estimate, and the limit
$(K_n-W_n)/n^{2}\to C_0$.  That last limit is the content of the source's
Lemmas~1 and~2, and the elementary proof of
Section~\ref{long1049:sec:cyclotomic-limit} replaces them, so they are not part of the
citation set.  Kernel-checked are the finite height comparisons of
Theorem~\ref{long1049:res:31over4}, the rational bracket of
Section~\ref{long1049:sec:regionbracket}, and the four integer identities of the degree
step in Section~\ref{long1049:sec:degrees}.  Checked by exact computation are the
reconstructions of $U_n$ and $V_n$ for $n\le4$, together with the numerical
values of $J$, $C_0$ and $\theta^{*}$; those are computational receipts, and
the reconstruction is described in Section~\ref{long1049:sec:receipts}.  No part of
Theorem~\ref{long1049:res:region} is a Lean theorem.

\paragraph{Attribution.}
Bundschuh and V\"a\"an\"anen proved the irrationality of this value on the
region $\log b/\log a<1/2-1/\pi^{2}=0.3986788163576622\ldots$
\cite[Thm.~2, p.~177; hypotheses pp.~175--176]{bv1994}, and every base of
Theorem~\ref{long1049:res:region} with $b\le3$ is already theirs.  Their printed
hypothesis for $\alpha=-1$ is $\lambda<(1/2+1/\pi^{2})^{-1}$ with
$\lambda=\log h(q)/\log|q|$; at $q=a/b$ this reads
$\log a/\log(a/b)<(1/2+1/\pi^{2})^{-1}$, which is the displayed inequality.  Duverney proved an
independent and strictly smaller region for the same series
\cite[Th\'eor\`eme~2, p.~174]{duverney1996}.  Zudilin announced a rational-base
extension of his generalized $q$-logarithm results under an assumption
$\log|r|>c\log|s|$ for a computable $c>0$, without computing a value
\cite{zudilin2016}.  Zudilin's 2004 paper supplies the forms and the exponent
$\mu$ under the standing hypothesis $p=1/q\in\Z\mathbin{\backslash}\{0,\pm1\}$, and it
states no rational-base result \cite[Sec.~2, p.~154; Thm.~1]{zudilin2004}.  The
contribution here is the rational specialisation of the 2004 forms, with the
denominator accounting and the limit passage carried out in full, which
identifies the printed $\mu$ as an admissible $c$ for $F$ itself, together with
the region that constant defines and its first new base.

\paragraph{Where the earlier criterion stops.}
Both regions are cut out by the same quantity.  The published cutoff
$1/2-1/\pi^{2}$ of \cite[Thm.~2, p.~177]{bv1994} is the reciprocal of
$\mu_{\mathrm{BV}}=2\pi^{2}/(\pi^{2}-2)$, the constant Van Assche later
recovered as an integer-base irrationality-exponent bound for $F(p)$
\cite[Thm.~1, p.~10]{vanassche2001}, and the cutoff $\theta^{*}$ of
Theorem~\ref{long1049:res:region} is the reciprocal of the smaller integer-base bound
$\mu=C_1/C_0$ printed at \cite[p.~162]{zudilin2004}.  The remark following
Theorem~\ref{long1049:res:archcap} records that reciprocal relation for every family
satisfying the hypotheses of that theorem, and Zudilin's family satisfies them
with decay exponent $\sigma=C_0$ and exact degree limit $d=C_1-C_0$.  The step
taken here is the homogenisation of Step~3 in the proof of
Theorem~\ref{long1049:res:region}, applied to the forms that carry the smaller exponent.
It substitutes one integer-base exponent bound for another, and it leaves the
analytic content of each bound where its own source proves it.  The bases
gained are exactly the strip $s^{\mu}<r\le s^{\mu_{\mathrm{BV}}}$, which is empty
for $s=2$ and $s=3$ and is first occupied at $s=4$ by the single numerator
$31$.  The kernel-checked declarations of Theorem~\ref{long1049:res:31over4} and of
Section~\ref{long1049:sec:regionbracket} decide the integer comparisons and the rational
bracket that place a base inside or outside each region.  They decide no step
of either analytic argument, and they improve neither exponent bound.

\paragraph{The exact remaining obligation.}
Theorem~\ref{long1049:res:region} settles no base with $\log b/\log a\ge\theta^{*}$, and
in particular settles nothing at $3/2$, whose parameter is
$\log2/\log3=0.6309297535714574\ldots$, a gap of $0.2252467\ldots$ above
$\theta^{*}$.  Enlarging the region by this proof requires a smaller integer-base
exponent bound carried by forms with the $\Z[p]$-integrality of Lemma~7 type and
with base-uniform size; a bound $2.4234$ would give $0.4126$ and would admit
$29/4$.  Formalising the theorem requires Lemma~7 and its inputs, the
identity~(9)--(11), and the steps of
Sections~\ref{long1049:sec:source-forms}--\ref{long1049:sec:integer-forms}.  Of these only the
integer degree bookkeeping is formalised here.

\begin{theorem}[degree-budget cap on the mechanism]\label{long1049:res:archcap}
Let $(U_n,V_n)\in\Z[x]^{2}$ be a sequence such that, for constants
$\sigma,\delta>0$ and $h\ge0$ independent of $n$ and of the base,
\begin{enumerate}
\item $\Lambda_n(x):=U_n(x)F(x)-V_n(x)\ne0$ for every real $x>1$;
\item $\deg U_n,\deg V_n\le\delta n^{2}(1+o(1))$;
\item $\log\max\bigl(H(U_n),H(V_n)\bigr)\le hn^{2}(1+o(1))$, where $H(P)$ is the
largest absolute value of a coefficient of $P$;
\item $\log|\Lambda_n(x)|=-\sigma n^{2}\log x\,(1+o(1))$ for every real $x>1$.
\end{enumerate}
Put $d_n:=\max(\deg U_n,\deg V_n)$.  Then $\sigma\le\delta$, and for every fixed
rational base $a/b>1$,
\[
 \limsup_{n\to\infty}n^{-2}\log\bigl|b^{d_n}\Lambda_n(a/b)\bigr|
 \le\delta\log b-\sigma\log(a/b).
\]
Consequently the homogenised forms tend to zero whenever
$\log b/\log a<\sigma/(\sigma+\delta)$, a sufficient region whose cutoff is at
most $1/2$.  If the actual degrees satisfy $d_n/n^{2}\to d$, then $d\ge\sigma$
and the limit exists and equals $d\log b-\sigma\log(a/b)$; in that case the
forms tend to zero below $\log b/\log a=\sigma/(\sigma+d)$ and their absolute
values tend to infinity above it, so the exact-degree case has no decaying
homogenised forms at $3/2$.
\end{theorem}

\begin{proof}
The displayed inequality is the identity
$n^{-2}\log|b^{d_n}\Lambda_n(a/b)|=(d_n/n^{2})\log b-\sigma\log(a/b)+o(1)$
together with hypothesis~(2); the degree bound is what makes
$b^{d_n}U_n(a/b)$ and $b^{d_n}V_n(a/b)$ integers.

For $\sigma\le\delta$, fix an integer $p\ge2$, write $\xi=F(p)$,
$\alpha=h+\delta\log p$ and $\beta=\sigma\log p$, and set
$\varepsilon_n=\Lambda_n(p)$.  Hypotheses~(2) and~(3) give the upper bound
$|U_n(p)|\le e^{(\alpha+o(1))n^{2}}$, an inequality and not an equality: a
degree and height bound leaves cancellation at $p$ possible.  Hypothesis~(4)
gives $\log|\varepsilon_n|=-\beta n^{2}(1+o(1))$ on both sides, and
hypothesis~(1) gives $\varepsilon_n\ne0$.  Since $U_n(p)\xi-V_n(p)=\varepsilon_n$
is a nonzero real tending to $0$ while $U_n(p)$ and $V_n(p)$ are integers, $\xi$
is irrational; in particular $U_n(p)\ne0$ for all large $n$, since otherwise
$\varepsilon_n=-V_n(p)$ would be a nonzero integer tending to $0$.

Let $P/Q$ be rational with $Q$ large, and let $n$ be least with
$|\varepsilon_n|<1/(2Q)$.  Minimality and the two-sided asymptotic give
$n^{2}\le(1+o(1))\log Q/\beta$, hence
\[
 |\varepsilon_n|\ge Q^{-1-o(1)},\qquad |U_n(p)|\le Q^{\alpha/\beta+o(1)} .
\]
If $U_n(p)P-V_n(p)Q\ne0$, then from
$U_n(p)P-V_n(p)Q=Q\varepsilon_n-U_n(p)(Q\xi-P)$ and $|Q\varepsilon_n|<1/2$ one
gets $|U_n(p)|\,|Q\xi-P|>1/2$, so
$|\xi-P/Q|>1/\bigl(2Q|U_n(p)|\bigr)$.  If instead $P/Q=V_n(p)/U_n(p)$, then
$|\xi-P/Q|=|\varepsilon_n|/|U_n(p)|$.  In both cases
$|\xi-P/Q|\ge Q^{-1-\alpha/\beta-o(1)}$, so the irrationality exponent satisfies
$\mu(\xi)\le1+\alpha/\beta$.  Dirichlet's theorem gives $\mu(\xi)\ge2$ for the
irrational $\xi$, hence $\beta\le\alpha$, that is
$(\sigma-\delta)\log p\le h$.  The three constants do not depend on $p$, so
letting $p\to\infty$ gives $\sigma\le\delta$.  The exact-degree statement is the
same argument with $d+\varepsilon$ in place of $\delta$ for every
$\varepsilon>0$.
\end{proof}

\begin{corollary}[Undivided forms below the square boundary]
\label{long1049:cor:no-decay-below-square}
Under the hypotheses of Theorem~\ref{long1049:res:archcap}, for positive integers
$a,b$ with $b<a<b^2$, the undivided forms $b^{d_n}\Lambda_n(a/b)$ do not
tend to zero.  No limit of $d_n/n^2$ is assumed.
\end{corollary}
\begin{proof}
Write $x=a/b$ and suppose $C_n=b^{d_n}\Lambda_n(x)\to0$.
Eventual nonvanishing gives
$\log|C_n|=d_n\log b+\log|\Lambda_n(x)|$ for all sufficiently large $n$.
Also $|C_n|\le1$ eventually.  The lower side of the remainder asymptotic
therefore implies
\[
 d_n\log b\le-\log|\Lambda_n(x)|
       =\sigma\log x\,n^2+o(n^2).
\]
Set $c=\sigma\log x/\log b$.  Since $1<x<b$, we have $0<c<\sigma$,
and the displayed inequality gives $d_n\le(c+\varepsilon)n^2$ eventually
for each $\varepsilon>0$.  The same family thus satisfies the cap hypotheses
with degree upper rate $c$, its original height bound, and its original
nonvanishing and remainder asymptotics at every real base greater than one.
The cap yields $\sigma\le c$, a contradiction.
\end{proof}
\noindent The no-decay conclusion is \href{https://github.com/wcook04/plectis-erdos/blob/3be82b1a7340284aea72e9a5c8493cb020843921/ErdosProblems/Erdos1049/PaperNoDecayR9.lean#L69}{kernel-checked in Lean} under the stated all-base hypotheses.


This strengthens the failure of a sufficient-cutoff test to an exclusion
of decay under the stated all-base hypotheses.  It does not supply an
actual-degree limit or assert divergence.  If $d_n/n^2\to d$, the separate
exact-degree result gives the stronger conclusion
$|b^{d_n}\Lambda_n(a/b)|\to\infty$ in the same strict region.
Equality $a=b^2$ remains unclassified.  The corollary concerns the undivided
forms; it does not exclude base-dependent content division or other
irrationality methods outside its hypotheses.

For the maximum-coefficient height convention, the passage to the
$\ell^1$ norm uses $\|P\|_1\le(\deg P+1)H(P)$.  The original quadratic
degree bound makes the extra logarithm $o(n^2)$.  This established height
bound remains available when the degree rate is replaced by $c$.



\noindent Hypothesis~(3) is essential and was absent from an earlier form of
this statement: a bound on $\deg U_n$ alone controls $|U_n(p)|$ only through
$H(U_n)$.  At the boundary $\log b/\log a=\sigma/(\sigma+d)$ the normalised
logarithm is zero and these hypotheses decide neither behaviour.
Theorem~\ref{long1049:res:archcap} constrains families satisfying its hypotheses and
does not exclude every possible Pad\'e construction.

For the family of Section~\ref{long1049:sec:source-forms} the fourth hypothesis is the
size estimate proved there and the second is~\eqref{long1049:eq:exact-degree}.  The third
is proved next, so the cap applies to that family with no further assumption.

\begin{lemma}[uniform coefficient height of the source
forms]\label{long1049:res:sourceheight}
There is a constant $h$ with
$\log\max\bigl(H(U_n),H(V_n)\bigr)\le hn^{2}$ for every $n\ge1$, where $U_n$ and
$V_n$ are the polynomials of~\eqref{long1049:eq:integer-polynomial-pair}.
\end{lemma}

\begin{proof}
Write $\|P\|$ for the sum of the absolute values of the coefficients of
$P\in\Z[X]$, so that $H(P)\le\|P\|$ and $\|PQ\|\le\|P\|\,\|Q\|$.

By Lemma~\ref{long1049:res:omega-indicator} every $\nu_{n,l}$ is $0$ or $1$, so
$D_N/\Omega_n=\prod_{l\in S}\Phi_l$ over a subset $S\subseteq\{1,\dots,N\}$.
Each $\Phi_l$ has Mahler measure $1$, and a polynomial of degree $D$ satisfies
$\|P\|\le2^{D}M(P)$, so $\|\Phi_l\|\le2^{\varphi(l)}$ and
\[
 \Bigl\|\frac{D_N}{\Omega_n}\Bigr\|\le2^{\sum_{l\le N}\varphi(l)}\le2^{225n^{2}} .
\]
For $A_n$, the Gaussian binomial $\genfrac{[}{]}{0pt}{}{m}{r}_X$ has nonnegative
coefficients summing to $\binom{m}{r}$, so
$\|c_{n,k}\|\le\binom{k-1}{a_1-1}\binom{\beta_n-a_2-1}{\beta_n-k-1}\le2^{2\beta_n}$
by~\eqref{long1049:eq:c-explicit}, and summing the at most $\beta_n$ terms
of~\eqref{long1049:eq:A-explicit} gives $\|A_n\|\le\beta_n2^{2\beta_n}$.  Hence
$\|U_n\|\le\|D_N/\Omega_n\|\,\|A_n\|\le e^{O(n^{2})}$ and
$H(U_n)\le e^{O(n^{2})}$.

For $V_n$, bound it on the circle $|z|=2$ and use Cauchy's estimate
$H(V_n)\le\max_{|z|=2}|V_n(z)|$, valid because $\deg V_n\ge0$ and each
coefficient is $|v_i|\le2^{-i}\max_{|z|=2}|V_n(z)|$.  On that circle
$|z^{l}-1|\ge2^{l}-1\ge1$, so each of the at most $\beta_n+a_0$ inner terms
of~\eqref{long1049:eq:B-explicit} has modulus at most $1$; also
$|c_{n,k}(z)z^{a_0k}|\le\|c_{n,k}\|\,2^{K_n}$ and
$|D_N(z)/\Omega_n(z)|\le\|D_N/\Omega_n\|\,2^{225n^{2}}\le e^{O(n^{2})}$, while
$|z^{-M_n}|\le1$.  Multiplying the $O(n)$ bounds and using $K_n=O(n^{2})$ gives
$\max_{|z|=2}|V_n(z)|\le e^{O(n^{2})}$, hence $H(V_n)\le e^{O(n^{2})}$.
\end{proof}

\noindent With Lemma~\ref{long1049:res:sourceheight}, Zudilin's family satisfies all four
hypotheses of Theorem~\ref{long1049:res:archcap}, with exact degree limit $d=C_1-C_0$
by~\eqref{long1049:eq:degree-limits}, and decay exponent $\sigma=C_0$ because
$\log\Lambda_n(x)=-(K_n-W_n)\log x+o(n^{2})$ for each fixed real $x>1$ by the
size estimate of Section~\ref{long1049:sec:integer-forms}.  For it $\sigma/(\sigma+d)=\theta^{*}$ and
$(\sigma+d)/\sigma=\mu$: within this family the rational-base threshold is the
reciprocal of the integer-base irrationality-exponent bound.
Lemma~\ref{long1049:res:sourceheight} is an ordinary proof and is not kernel-checked.

\paragraph{The region, explicitly.}
Among coprime $a/b$ with $a\le60$ the region of Theorem~\ref{long1049:res:region} has
$137$ members, of which $78$ are non-integral.  The tightest members are $53/5$,
at margin $0.000313$ below $\theta^{*}$, and $31/4$, at margin $0.001985$; the
closest miss is $52/5$ at $\theta=0.4073243836\ldots$, short by $0.001641$.  The
bases new relative to \cite{bv1994} are those in the strip
$s^{\mu}<r\le s^{\mu_{\mathrm{BV}}}$, which is empty for $s=2$ and $s=3$, is
$\{31\}$ for $s=4$, and is $\{53,54,56\}$ for $s=5$.  The strip is infinite:
its width $s^{\mu_{\mathrm{BV}}}-s^{\mu}$ eventually exceeds $s$, so for every
large enough denominator it contains an integer congruent to $1$ modulo $s$.
Hence $31/4$ is the new base of least denominator and least numerator.

\paragraph{The direction is optimal in its box.}
Over the $37{,}533$ primitive directions on Zudilin's cone with all entries at
most $30$, the quantity $\theta^{*}(\mathrm{dir})=C_0/C_1$ computed from his
(25) and (26) is maximised by $(14,12,14;27)$ and its group image
$(15,12,13;26)$, both at $0.4056830213840605\ldots$; the next value is
$0.4056394327738419\ldots$, at $(16,13,14;28)$.  Larger boxes were not scanned,
so this is a search result over that box and not a proof of optimality over the
whole cone.  The enumeration ranges over integer tuples
$(\alpha_0,\alpha_1,\alpha_2;\beta)$ with $1\le\alpha_j\le30$,
$1\le\beta\le30$, satisfying the cone conditions of \cite[Sec.~5]{zudilin2004}
and taken up to the common factor of the four entries; the objective is the
ratio $C_0/C_1$ of the source's (25) and (26), evaluated by the same
thirteen-interval trigamma sum.  The two maximising tuples are exchanged by the
source's permutation group, so they are one direction.

\subsection{Finite receipts}\label{long1049:sec:receipts}

Four finite computations support statements above.  Each is exact rational or
integer arithmetic except where a numerical enclosure is named; none of them
proves a statement quantified over all indices, and none is kernel-checked.

The forms of Section~\ref{long1049:sec:source-forms} were reconstructed as explicit
elements of $\Z[X]$ from cyclotomic products, without rational-function
arithmetic, for $n\le4$.  At $n=1,2,3$ the reconstruction confirms
$\deg A_n=K_n$ with $K_n=587,2264,5032$, that $\Omega_n$ divides both $D_NA_n$
and $D_NB_n$ with zero remainder, that $X^{M_n}$ divides both quotients, and
that $W_n=\deg U_n=333,1315,2944$ with $\deg V_n=W_n-1$ and leading
coefficients $\pm1$.  It also confirms
$\sum_{l\le N}\varphi(l)=72,278,628$ and $\Sigma_n=25,94,219$, hence
$(K_n-W_n)/n^{2}=254,949/4,232$.  The identity
$U_n(x)F(x)-V_n(x)=x^{-M_n}(D_N/\Omega_n)(x)H_n(x)$ was then evaluated at
$x=31/4$, $3$ and $7/2$ to relative accuracy below $10^{-39}$ at $n=3$, with
$\widehat\Lambda_n>0$ at each, with $b^{W_n}U_n(a/b)$ an integer and
$b^{W_n-1}U_n(a/b)$ not an integer at $31/4$ and $7/2$, and with $H_n(x)$
inside the bounds of Section~\ref{long1049:sec:integer-forms}.

The constant was evaluated independently of the source's printed digits from
$\omega$, the thirteen intervals of Lemma~\ref{long1049:res:omega-indicator} and the
trigamma series at forty digits, giving
$J=77.94318447500909\ldots$, $C_0=221.30008816500502\ldots$ and
$C_1/C_0=2.46497868357497\ldots$, which agree with
\cite[p.~162]{zudilin2004} to its printed precision.

The main term $K_n\log b-(K_n-W_n)\log a$ at $31/4$, computed from
$\sum_{l\le N}\varphi(l)$ and $\Sigma_n$ alone, is negative from $n=1$ and its
normalisation by $n^{2}$ runs $-58.478,-10.280,-4.297,-3.863$ at
$n=1,10,100,400$ against the limit $C_1\log4-C_0\log31=-3.718\ldots$.  The
elementary bound on the convergence rate is $O(\log^2(n+2)/n)$.
Indeed, summing the $O(y\log(2y))$ endpoint errors over the floor blocks
with $k\le n$ gives $O(n\log^2(n+2))$; the remaining main-term tail is
$O(1)$ since its summands are $O(n^2/k^3)$.  The linear terms of $M_n$
add $O(n)$.  The stronger rate $O(1/n)$ does not follow from these estimates.  Only the limit enters Theorem~\ref{long1049:res:region}.

The order and leading coefficient of $V_N^{*}$ were computed exactly for
$1\le N\le7$, giving orders $0,1,5,14,30,55,91$ and coefficients
$1,6,108,4320,324000,40824000,8001504000$.

\section{The sharp $41/65$ certificate}\label{long1049:sec:sharp}

The numerical obstruction at $3/2$ is controlled by one small exact
calculation.  Writing it out is useful because the upper and lower inequalities
play different roles: the upper inequality proves every strict analytic
threshold below, while the lower inequality establishes the sharp rational
scale.  Exact failure of the rank-$41$ selector count one row earlier uses the
additional integer comparison $2^{129}<3^{82}$ below.

\begin{theorem}[sharp power bracket]\label{long1049:res:powerbracket}
One has
\[
 2^{64}<3^{41}<2^{65}.
\]
Consequently
\[
 \frac{41}{65}<\frac{\log2}{\log3},
 \qquad
 \frac{\log3}{\log2}<\frac{65}{41}.
\]
\end{theorem}

\begin{proof}
Direct evaluation gives
\[
 \begin{aligned}
 2^{64}&=18446744073709551616,\\
 3^{41}&=36472996377170786403,\\
 2^{65}&=36893488147419103232.
 \end{aligned}
\]
The logarithmic inequalities follow by taking logarithms and dividing by the
positive numbers $41\log3$ and $41\log2$, respectively.
\end{proof}

\noindent The integer sides are checked as
\lword{Erdos1049/AdelicHeightBridge.lean}{61}{threePow_fortyOne_lt_twoPow_sixtyFive}{the upper certificate}
and
\lword{Erdos1049/AdelicHeightBridge.lean}{64}{twoPow_sixtyFour_lt_threePow_fortyOne}{the sharp lower certificate}.

The upper half is stronger than the earlier $81/200$ comparison in the
direction needed at $3/2$.  Combining it with the elementary bound
$1/2-1/\pi^2<2/5$ gives the following exact deficits.

For $\rho,\sigma\in\R$, write
\[
 \Theta_{\mathrm{HP}}(\rho,\sigma)=
 \frac{(1+\rho^2)/2+\sigma-3\sigma^2/\pi^2}
 {(1+\rho)^2/2+\sigma(1+\rho)+(1+\rho^2)/2+\sigma}.
\]

\begin{corollary}[height and Hankel deficits]\label{long1049:res:sharpgaps}
For every $\rho,\sigma\in\R$ with $0\le\rho$ and $1+\rho\le\sigma$,
\[
 \frac3{13}<\frac{\log2}{\log3}
   -\left(\frac12-\frac1{\pi^2}\right),
 \qquad
 \frac3{13}<\frac{\log2}{\log3}-\Theta_{\mathrm{HP}}(\rho,\sigma),
\]
where $\Theta_{\mathrm{HP}}$ is the rectangular exponent threshold.  Moreover
\[
 \frac{\log3/\log2-1}{3}<\frac8{41}.
\]
\end{corollary}

\begin{proof}
The first inequality follows from
$41/65-2/5=3/13$, Theorem~\ref{long1049:res:powerbracket}, and
$1/2-1/\pi^2<2/5$.  The rectangular threshold is no larger than the classical
margin $1/2-1/\pi^2$, so the second follows.  The final inequality is a direct
rearrangement of $\log3/\log2<65/41$.
\end{proof}

\noindent The uniform rectangular bound and cubic charge bound are
\lword{Erdos1049/AdelicHeightBridge.lean}{103}{threeHalves_rectangular_hp_gap_gt_threeThirteenths}{checked here}
and
\lword{Erdos1049/AdelicHeightBridge.lean}{114}{threeHalves_hankelChargeThreshold_lt_eightFortyOne}{checked here}.

The next statement is the source-facing charge comparison.  It does not prove
the degree ceilings: those come from the scalar factor and the universally
forced first-order southeast-border factor in the Zudilin model.  It proves
that even granting those ceilings, neither extraction reaches the required
$39/41$ fraction of the raw charge.

\begin{theorem}[scalar and border charge no-go]\label{long1049:res:chargeceilings}
For every integer $N>0$,
\[
 41(N^3-N)<39(4N^3-3N^2).
\]
For every integer $N\ge2$,
\[
 41(2N^3-N)<39(4N^3-3N^2).
\]
Hence the same strict inequalities hold with the left side replaced by
$41E$ whenever, respectively, $E\le N^3-N$ or $E\le2N^3-N$.
\end{theorem}

\begin{proof}
After subtraction, the first inequality is
\[
 N(115N^2-117N+41)>0,
\]
which holds for $N>0$.  The second becomes
\[
 N(74N^2-117N+41)>0.
\]
For $N\ge2$ the quadratic factor is positive and increasing.  The assertions
for $E$ follow by monotonicity.
\end{proof}

\noindent The application-facing downward-closed forms are
\lword{Erdos1049/AdelicHeightBridge.lean}{140}{zudilinScalarContent_cannot_meet_required_charge}{the scalar no-go}
and
\lword{Erdos1049/AdelicHeightBridge.lean}{170}{zudilinScalarPlusBorder_cannot_meet_required_charge}{the scalar-plus-border no-go}.

The formal source checks Theorem~\ref{long1049:res:powerbracket},
Corollary~\ref{long1049:res:sharpgaps}, and Theorem~\ref{long1049:res:chargeceilings} as exact
Lean propositions.  Their role is exclusion: scalar content and the forced
first-order border do not supply enough charge.  Higher residual valuations,
determinant cancellation, a different integral model, and irrationality of
$F(3/2)$ remain open.

\subsection{The sharp $q$-order of the normalised Hankel
determinant}\label{long1049:sec:hankel-order}

The charge comparison of Corollary~\ref{long1049:res:sharpgaps} is a statement about how
much denominator the Hankel determinant can shed.  How much it actually sheds
is settled by the following exact order.

\begin{theorem}[sharp all-rank $q$-order and leading
coefficient]\label{long1049:res:zudilin-sharp-qorder}
For every rank $N$, the normalised Hankel determinant $V_N^{*}$ of
\cite[Sec.~4]{zudilin2016}, evaluated at $x=z=1$, has
\[
 \operatorname{ord}_q V_N^{*}=\frac{N(N-1)(2N-1)}{6},
 \qquad
 \text{leading coefficient}\quad\frac{(N!)^{2}(N+1)!}{2^{N}} .
\]
\end{theorem}

\noindent The source proves the inequality
$\operatorname{ord}_q V_N^{*}\ge N(N-1)(2N-1)/6$ and uses it as an upper bound
on $|V_N^{*}|$ \cite[Sec.~4]{zudilin2016}.  Equality is the statement that no
further cancellation is available inside the determinant, so the hope of a
hidden cubic decay there is closed.

\paragraph{The row identity.}
Work over $A=\Z[[q]]$.  Let $H(X)=1+\sum_{s\ge1}a_s(q)X^{s}$ and put
\[
 W_m(t)=q^{(m+1)t}\prod_{r=1}^{m}H(q^{r}),
 \qquad
 D_j=\prod_{r=0}^{j-1}(I-q^{r}\mathcal N),
 \qquad
 E(m,j)=mj-\frac{j(j-1)}2,
\]
where $\mathcal N$ is the backward shift $(\mathcal Nf)_m=f_{m-1}$ in the index
$m$, so $D_j$ is the source's backward-difference operator of
\cite[Sec.~4]{zudilin2016}.  Write $\bar H=H\bmod q$ and
$h_r=[X^{r}]\bar H(X)^{-1}$, with $h_r=0$ for $r<0$ and $h_0=1$.

\begin{lemma}[all-depth initial monomial]\label{long1049:res:allrowinitial}
For $m\ge j\ge0$ one has $D_jW_m(t)\in q^{E(m,j)}A$ and
\[
 \bigl[q^{E(m,j)}\bigr]D_jW_m(t)=(-1)^{j}h_{j-t}.
\]
\end{lemma}

\begin{proof}
Define the following transition coefficients on states $(e_u)_{u\ge0}$.
Powers $\Phi^j$ mean the finite-depth path sums, not endomorphisms of the
algebraic direct sum: for fixed initial state $t$, endpoint $u$ and depth $j$,
all visited states are at most $\max(t,u+j)$, so each such coefficient is a
finite sum.  Subsequent sums over endpoints are interpreted coefficientwise
as justified below.  The transition coefficients are
\[
 \Phi e_0=\sum_{s\ge1}a_se_{s-1},
 \qquad
 \Phi e_u=(q^{u}-1)e_{u-1}+q^{u}\sum_{s\ge1}a_se_{u+s-1}\quad(u>0).
\]
A direct expansion gives
$W_m(u)-W_{m-1}(u)=q^{m}\sum_v(\Phi e_u)_vW_{m-1}(v)$: both sides equal
$q^{mu}\prod_{r<m}H(q^{r})\bigl(q^{u}H(q^{m})-1\bigr)$.  Using
$r+E(m-1,r)=E(m,r)$, induction on $j$ turns this into
\begin{equation}\label{long1049:eq:state-expansion}
 D_jW_m(t)=q^{E(m,j)}\sum_u(\Phi^{j}e_t)_uW_{m-j}(u).
\end{equation}
The sum is locally finite: one transition lowers a state by at most one, so a
depth-$j$ path ending at $u$ never visits a state above $\max(t,u+j)$, and
$\operatorname{ord}W_{m-j}(u)=(m-j+1)u$, so only finitely many endpoints
contribute to any fixed power of $q$.

Modulo $q$ the operator simplifies: $\bar\Phi e_u=-e_{u-1}$ for $u>0$, and
$\bar\Phi e_0=\sum_{s\ge1}a_s(0)e_{s-1}$.  Put $c_{j,t}=(\bar\Phi^{j}e_t)_0$.
Then $c_{j+1,t}=-c_{j,t-1}$ for $t>0$ and
$c_{j+1,0}=\sum_{s=1}^{j+1}a_s(0)c_{j,s-1}$, the sum terminating because a state
above $j$ cannot reach $0$ in $j$ steps.  Now $c_{0,t}=\delta_{t,0}=h_{-t}$, and
if $c_{j,t}=(-1)^{j}h_{j-t}$ for all $t$ then $c_{j+1,t}=(-1)^{j+1}h_{j+1-t}$ for
$t>0$ at once, while for $t=0$ the identity $\bar H\bar H^{-1}=1$ gives
$h_{j+1}=-\sum_{s\ge1}a_s(0)h_{j+1-s}$ and hence
$c_{j+1,0}=(-1)^{j}\sum_{s\ge1}a_s(0)h_{j+1-s}=(-1)^{j+1}h_{j+1}$.  So
$c_{j,t}=(-1)^{j}h_{j-t}$ for all $j,t$.  Since
$\operatorname{ord}W_{m-j}(u)=(m-j+1)u$ and $m\ge j$, in
\eqref{long1049:eq:state-expansion} only the state $u=0$ contributes at degree $E(m,j)$,
which gives both assertions.
\end{proof}

\noindent Apply this to the normalised source tail
\[
 T_{m,t}=q^{(m+1)t}\,(q;q)_m^{3}\,\frac{(q^{t+1};q)_m}{(q^{m+t+1};q)_{m+1}} .
\]
With $C_t=(1-q^{t+1})^{-1}$ and
\[
 H_t(X)=\frac{(1-X)^{3}(1-q^{t}X)^{2}}{(1-q^{t}X^{2})(1-q^{t+1}X^{2})}
\]
one has $T_{m,t}=C_tW_m^{H_t}(t)$, because
$\prod_{r\le m}H_t(q^{r})$ telescopes to
$(q;q)_m^{3}(1-q^{t+1})(q^{t+1};q)_m/(q^{m+t+1};q)_{m+1}$: the two denominator
products contribute the consecutive factors $1-q^{t+2},\dots,1-q^{t+2m+1}$.
Reducing modulo $q$ gives $\bar H_t=(1-X)^{3}$ for $t>0$, so
$h^{(t)}_r=\binom{r+2}{2}$, and $\bar H_0=(1-X)^{4}/(1+X)$, so
$h^{(0)}_r=(r+1)(r+2)(2r+3)/6$.  Both have $C_t(0)=1$.  The rows are
$v_m=\sum_{t\ge0}T_{m,t}$, and interchanging that sum with the finite
difference $D_j$ is legitimate because $\operatorname{ord}T_{m,t}=(m+1)t$.  Terms
with $t>j$ contribute $h^{(t)}_{j-t}=0$, so Lemma~\ref{long1049:res:allrowinitial} at
$m=j+\ell$, where $E(j+\ell,j)=j(j+1)/2+j\ell$, gives
\begin{equation}\label{long1049:eq:row-initial}
 D_jv_{j+\ell}=(-1)^{j}\frac{(j+1)^{2}(j+2)}2\,q^{\,j(j+1)/2+j\ell}
 +O\bigl(q^{\,j(j+1)/2+j\ell+1}\bigr)
 \qquad(j,\ell\ge0),
\end{equation}
the coefficient because
\[
 h^{(0)}_j+\sum_{t=1}^{j}h^{(t)}_{j-t}
 =\frac{(j+1)(j+2)(2j+3)}6+\binom{j+2}3
 =\frac{(j+1)^{2}(j+2)}2 .
\]

\begin{proof}[Proof of Theorem~\ref{long1049:res:zudilin-sharp-qorder}]
The operators $D_j$ act by lower unitriangular row operations, so they leave
$\det(v_{i+j})_{0\le i,j<N}$ unchanged.  By~\eqref{long1049:eq:row-initial} the entry in
row $j$ and column $\ell$ has order $e(j,\ell)=j(j+1)/2+j\ell$.  In the Leibniz
expansion the weight of a permutation $\varsigma$ is
$\sum_j\bigl(j(j+1)/2+j\varsigma(j)\bigr)$, and by the rearrangement inequality
$\sum_jj\varsigma(j)$ is uniquely minimised by the reversal
$\varsigma(j)=N-1-j$, the values $j$ being distinct.  The minimum weight is
$\sum_{j<N}j^{2}=N(N-1)(2N-1)/6$, so exactly one Leibniz term attains it and no
cancellation is possible there.  The sign of the reversal is
$(-1)^{N(N-1)/2}$, which cancels $\prod_{j<N}(-1)^{j}$, and the surviving
coefficient is
\[
 \prod_{j=0}^{N-1}\frac{(j+1)^{2}(j+2)}2=\frac{(N!)^{2}(N+1)!}{2^{N}} .
\]
\end{proof}

\paragraph{Formal order and analytic size are different questions.}
Theorem~\ref{long1049:res:zudilin-sharp-qorder} fixes the first nonzero power of $q$ and
its coefficient at each fixed rank.  It does not control the value at a fixed
rational $q$ as the rank grows.  For the logic, the integer polynomials
$f_N(q)=C_Nq^{B_N}(1-q)^{N^{3}}$, with $B_N=\sum_{j<N}j^{2}$ and
$C_N=(N!)^{2}(N+1)!/2^{N}$, have exactly the order $B_N$ and exactly the leading
coefficient $C_N$, while $f_N(2/3)=C_N(2/3)^{B_N}3^{-N^{3}}$ carries a further
cubic exponential factor that neither datum sees.  The following separate
positive-measure argument supplies the fixed-base estimate for $V_N^*$;
it is not inferred from the formal order.

\paragraph{A separate positive-measure estimate.}
For fixed $0<q<1$ write $P=(q;q)_\infty$, $Q=(\sqrt q;q)_\infty$,
$T=(-1;q)_\infty^2$, and
\[
 G_q(w)=\frac1{(w;q)_\infty^3}\sum_{t\ge0}\frac{w^t}{(q;q)_t}
          \frac{(q^tw^2;q)_\infty}{(q^tw;q)_\infty^2},
 \qquad \gamma_k=[w^k]G_q(w),\qquad c_k=\frac{(k+1)^2(k+2)}2.
\]
All products converge normally on $|w|\le r<1$ and the $t$-sum is bounded
there by a constant times $\sum_t r^t$.  The finite-to-infinite product
identities give $v_m^*=P^4G_q(q^{m+1})$.
The identity
\[
 \frac{(Aw;q)_\infty}{(w;q)_\infty}
   =\sum_{j\ge0}\frac{(A;q)_j}{(q;q)_j}w^j\quad(0\le A\le1)
\]
follows by comparing coefficients in
$(1-w)R(w)=(1-Aw)R(qw)$ with $R(0)=1$; the series converges for $|w|<1$
since its coefficients are at most $P^{-1}$.  They are nonnegative, and
at least $(A;q)_\infty$ if $A<1$.  For $A=1$ the series equals $1$.
Factor the numerator of the $t$th term using
\[
 (q^tw^2;q)_\infty=(q^{t/2}w;q)_\infty(-q^{t/2}w;q)_\infty
 (q^{(t+1)/2}w;q)_\infty(-q^{(t+1)/2}w;q)_\infty.
\]
Pair the two positive-argument factors with two copies of
$(w;q)_\infty^{-1}$.  All remaining factors have nonnegative
coefficients.  The $t=0$ term bounds $\gamma_k$ below by
$Q\binom{k+3}3$.  For upper bounds, a nonnegative-coefficient regular
factor $R$ with $R(1)\le C$, multiplying a series with nondecreasing
coefficients $d_k$, contributes at most $Cd_k$ by convolution.
This rule bounds the $t=0$ term by $TP^{-4}\binom{k+3}3$ and the sum
of the $t\ge1$ terms by $TP^{-6}\binom{k+2}3$.  Consequently
\[
 \frac Q3 c_k\le\gamma_k\le T(P^{-4}+P^{-6})c_k.
\]
Thus $\sum_{k\ge0}P^4\gamma_kq^k\delta_{q^k}$ is a finite positive measure
with infinitely many distinct support points and moments $v_m^*$.
A nonzero polynomial of degree less than $N$ cannot vanish at all of its
first $N$ atoms; the associated Gram matrix is positive definite.
Finite Cauchy--Binet followed by convergence of every matrix entry and
monotone convergence of the nonnegative tuple sums gives
\[
 V_N^*=\sum_{k_0<\cdots<k_{N-1}}
 \prod_i(P^4\gamma_{k_i}q^{k_i})\prod_{i<j}(q^{k_i}-q^{k_j})^2.
\]
Retaining the tuple $k_i=i$ and using
$\prod_{d=1}^{N-1}(1-q^d)^{2(N-d)}\ge P^{2N}$ gives the lower bound
below.  For the upper bound put $k_i=i+\lambda_i$, use
$c_{i+\lambda}/c_i\le(\lambda+1)^3$, discard residual Vandermonde factors
bounded by $1$, and drop the ordering restriction on the nonnegative
$\lambda_i$.  With $S(q)=(1+4q+q^2)/(1-q)^4$ this gives
\[
 (P^6Q/3)^N C_Nq^{B_N}\le V_N^*(q)
 \le\bigl[P^4T(P^{-4}+P^{-6})S(q)\bigr]^N C_Nq^{B_N}.
\]
Both constants are positive and finite, so $V_N^*(q)>0$ and
$\log(V_N^*(q)/(C_Nq^{B_N}))=O_q(N)$, including $N=0$ under the empty
determinant convention.  This is a separate analytic proof; it supplies no
new denominator factor for the 2004 polynomial forms.
The generating-function identification and infinite determinant expansion
in this argument are not yet formalised in Lean.

\paragraph{Evidence decomposition.}
The proof above is ordinary mathematics.  Three parts of it are kernel-checked
in this library.  The all-depth
associated-graded recurrence is
\lword{Erdos1049/AdelicHeightBridge.lean}{200}{hankelAssociatedCoeff_eq_reciprocal}{checked},
the initial monomial of transformed row $1$ in every column, of order exactly
$l+1$ and coefficient exactly $-6$, is
\lword{Erdos1049/AdelicHeightBridge.lean}{1422}{zudilinTransformedNormalizedMoment_one_initialMonomial}{checked},
and the two closed forms $6\operatorname{ord}=N(N-1)(2N-1)$ and
$2^{N}\mathrm{lc}=(N!)^{2}(N+1)!$ are assembled in
\lword{Erdos1049/AdelicHeightBridge.lean}{1753}{zudilinSharpHankelOrderAndCoeff_algebraicAssembly}{one
checked theorem}.  In the current public, kernel-checked library, row $j=2$,
with initial monomial $18q^{2l+3}$, is checked in a module of the
external-verification cut, while rows $j\ge3$ are not kernel verified.  A later
all-row Lean source candidate removes the row hypothesis at source level, but
its focused build and axiom audit are marked \textsc{unrun}, and it is not
included in this public release.  The Lean identification of $V_N^{*}$ itself
with the assembled leading matrix therefore stays conditional on a row
hypothesis here.  Independently of the proof, the order and the
leading coefficient were computed exactly for $1\le N\le7$, giving orders
$0,1,5,14,30,55,91$ and leading coefficients
$1,6,108,4320,324000,40824000,8001504000$, each matching the closed forms.

\paragraph{Attribution and remaining obligation.}
The antecedent is the inequality of \cite[Sec.~4]{zudilin2016}; the equality and
the leading coefficient are proved here.  Lemma~\ref{long1049:res:allrowinitial}
and~\eqref{long1049:eq:row-initial} give the initial monomial of every row.  The
later all-row Lean source is an unrun candidate rather than verification
authority, so the current public Lean statement about $V_N^{*}$ retains a row
hypothesis that the proof above discharges.  Nothing in this subsection decides
the arithmetic nature of $F(3/2)$.

\section{Integer scalar content is neutral}\label{long1049:sec:primitive}

An irrationality argument by linear forms replaces the clearing scheme by
explicit rational approximation: one constructs integer linear forms in $1$ and
$F(\beta)$ whose analytic decay outruns the height of their common denominator.
For an integer coefficient pair \((U,V)\) and a real target \(S\), put
\[
 L_S(U,V)=US-V,
 \qquad
 \Delta\bigl((U_n,V_n),(U_m,V_m)\bigr)=U_nV_m-U_mV_n.
\]
We call \(L_S(U,V)\) the \emph{error} of the row at \(S\); a good approximation
is one that makes it small.  The second expression is the exterior determinant
of the two rows, and it eliminates \(S\) exactly:
\[
\Delta=U_mL_S(U_n,V_n)-U_nL_S(U_m,V_m).
\]

This identity is what makes the determinant useful.  Since \(\Delta\) is an
integer, if it does not vanish then
\[
 1\le|\Delta|
  \le|U_m|\,\bigl|L_S(U_n,V_n)\bigr|+|U_n|\,\bigl|L_S(U_m,V_m)\bigr|,
\]
and the two errors cannot both be smaller than
\(1/(|U_n|+|U_m|)\).  The determinant therefore carries two competing
quantities at once, an integer that divides it and its own absolute value, and
the theorem below is about how a rescaling moves them.

\begin{theorem}[integer-scalar-content no-go]\label{long1049:res:content}
Let \(S\) be real, let \((U_n,V_n)\) and \((U_m,V_m)\) be pairs of integers, and
let \(c_n,c_m\) be integers.  Then
\[
 L_S(c_nU_n,c_nV_n)=c_nL_S(U_n,V_n),
\]
\[
 \Delta\bigl(c_n(U_n,V_n),c_m(U_m,V_m)\bigr)
 =c_nc_m\Delta\bigl((U_n,V_n),(U_m,V_m)\bigr),
\]
and consequently
\[
 \left|\Delta\bigl(c_n(U_n,V_n),c_m(U_m,V_m)\bigr)\right|
 =|c_n|\,|c_m|\,
   \left|\Delta\bigl((U_n,V_n),(U_m,V_m)\bigr)\right|.
\]
In particular \(c_nc_m\) divides the scaled determinant.  Hence a local
divisor supplied only by the two integer scalar factors is paid for by exactly
the same factor in the Archimedean determinant height.
\end{theorem}

\begin{proof}
All three identities are routine expansions in \(\Z\) or \(\R\); the
divisibility statement uses the primitive determinant as its witness.
\end{proof}

Informally, Theorem~\ref{long1049:res:content} says that multiplying specialised integer
rows by scalar factors moves the local gain and the Archimedean cost by exactly
the same amount.  Within an argument whose only extra divisor is integer scalar
content, primitive normalisation therefore loses no net gain.  This says
nothing about polynomial factors before specialisation, cross-row common
factors, determinant-specific arithmetic, or additive combinations.

\begin{example}\label{long1049:ex:content}
Take \((U_n,V_n)=(1,2)\) and \((U_m,V_m)=(3,5)\), so that
\(\Delta=1\cdot5-3\cdot2=-1\).  Multiplying the first row by \(c_n=6\) and the
second by \(c_m=10\) gives the rows \((6,12)\) and \((30,50)\), whose
determinant is \(6\cdot50-30\cdot12=-60\).  That determinant is now divisible
by \(60\), which looks like a local gain of \(60\); and its absolute value has
risen from \(1\) to \(60\), which is a cost of exactly the same size.
\end{example}

\noindent Lean checks the error identity in
\lword{Erdos1049/RationalPadeArithmetic.lean}{84}{rationalPadeError_mul}{error
scaling}, the determinant identity in
\lword{Erdos1049/RationalPadeArithmetic.lean}{94}{rationalPadeExteriorDet_mul_contents}{content
factorisation}, the exact absolute-height identity in
\lword{Erdos1049/RationalPadeArithmetic.lean}{104}{natAbs_rationalPadeExteriorDet_mul_contents}{absolute
determinant scaling}, and the divisor statement in
\lword{Erdos1049/RationalPadeArithmetic.lean}{116}{contentProduct_dvd_rationalPadeExteriorDet}{content-product
divisibility}.  The elimination identity is the checked
\lword{Erdos1049/RationalPadeArithmetic.lean}{124}{rationalPadeExteriorDet_cast_eq}{exterior
determinant identity}.

The key point is that the two scalings are the same scaling.  Each identity on
its own is a one-line expansion, and the interest of the theorem is not in any
one of them.  Taken together they say that the factor \(c_nc_m\) which a
rescaling introduces into the determinant reappears undiminished, as
\(|c_n|\,|c_m|\), in the absolute value of that determinant.  The third
identity is displayed with absolute values for exactly that reason: it is what
makes the statement one about the Archimedean height and not about divisibility
alone.  Whatever \(c_n\) and \(c_m\) are, a rescaling therefore leaves the
balance between the local divisor and that height where it was.

The theorem does not construct primitive Pad\'e rows, estimate their
remainders, or prove that their exterior determinant is nonzero.  It removes
one source of apparent gain: multiplying a useful row by a large common
integer cannot improve the local-to-Archimedean balance.  Within a construction
whose proposed gain comes only from those integer scalars, the rows may be
primitive-normalised without a net loss.  This does not say that every
candidate family must use such a normalisation or that every proof needs
separate \(2\)-adic and \(3\)-adic gain.  The theorem quantifies over arbitrary
integer pairs, so it applies whenever a construction has reached that
specialised-row stage.

\section{Endpoint residues at $(3,2)$ and the four-jet
kernel}\label{long1049:sec:endpoints}

Zudilin's treatment of \(q\)-harmonic series~\cite{zudilin2004} builds linear
forms of the shape used in Section~\ref{long1049:sec:primitive} out of Heine's basic
transform.  The generic endpoint and jet lemmas below concern arbitrary
integral coefficient pairs of this shape; they do not construct or instantiate
Zudilin's actual polynomial family.  Each lemma uses no property beyond
integrality and quantifies over arbitrary elements of \(\Z[X]\).  The exception is
Theorem~\ref{long1049:res:scalar} at the end of the section, whose subject is the scalar
parameters of Zudilin's cone.  It says nothing about the coefficient
polynomials.

Substituting \(X=3/2\) into an integer polynomial produces a rational number,
and multiplying by a power of \(2\) clears its denominator.  The following
evaluation records that cleared numerator, so that all the arithmetic below
stays inside \(\Z\).  For \(P(X)=\sum_i p_iX^i\in\Z[X]\) and a declared width
\(W\ge0\), put
\[
 H_W(P)=\sum_{i=0}^{W}p_i\,3^i2^{W-i}.
\]
This is the
\lword{Erdos1049/ZudilinConeArithmetic.lean}{98}{homEvalThreeTwo}{homogeneous
endpoint evaluation}.  It is homogeneous in the sense that \(X^{i}\) is
replaced by \(3^{i}2^{W-i}\), so the numerator and the denominator of the base
are carried symmetrically.  When \(W\ge\deg P\) it is exactly the cleared
numerator, since
\[
 \sum_{i=0}^{W}p_i\,3^i2^{W-i}=2^{W}\sum_{i=0}^{W}p_i\left(\tfrac32\right)^{i}
 =2^{W}P\!\left(\tfrac32\right).
\]

\begin{theorem}[endpoint residues]\label{long1049:res:endpoints}
Let \(P=\sum_ip_iX^i\in\Z[X]\) and let \(W\ge0\).  Then
\[
 H_W(P)\equiv p_0\,2^W\pmod 3,\qquad
 H_W(P)\equiv p_W\,3^W\pmod 2.
\]
Consequently a unit constant coefficient prevents divisibility by \(3\), and
a unit coefficient at the declared width \(W\) prevents divisibility by \(2\).
\end{theorem}

\begin{proof}
Modulo \(3\), every summand with \(i>0\) vanishes; modulo \(2\), every summand
with \(i<W\) vanishes.  The remaining powers are units in the corresponding
residue fields.
\end{proof}

Since \(2^{W}\) is invertible modulo \(3\) and \(3^{W}\) is invertible modulo
\(2\), the two congruences say more than the stated consequence: divisibility
of \(H_W(P)\) by \(3\) is decided by \(p_0\) alone, and divisibility by \(2\) by
\(p_W\) alone.  The rest of the coefficient vector is invisible to both primes.
The coefficient \(p_W\) is the top coefficient of \(P\) exactly when
\(\deg P=W\), and is zero when \(\deg P<W\); the identity
\(H_W(P)=2^{W}P(3/2)\) likewise holds only for \(\deg P\le W\).

\begin{example}\label{long1049:ex:endpoints}
Take \(W=2\).  The three polynomials below differ only at an endpoint.
\begin{center}
\begin{tabular}{@{}llcc@{}}
\toprule
\papertablehead
\(P\) & \(H_2(P)\) & \(3\mid H_2(P)\) & \(2\mid H_2(P)\)\\
\midrule
\(X^2+1\)  & \(1\cdot4+0\cdot6+1\cdot9=13\) & no  & no \\
\(X^2+3\)  & \(3\cdot4+0\cdot6+1\cdot9=21\) & yes & no \\
\(2X^2+1\) & \(1\cdot4+0\cdot6+2\cdot9=22\) & no  & yes \\
\bottomrule
\end{tabular}
\end{center}
The first has both endpoints equal to \(1\) and its evaluation, \(13\), is
divisible by neither prime; as a check, \(2^{2}\bigl((3/2)^2+1\bigr)=13\).  The
second and third show that each hypothesis is used: spoiling the constant
endpoint admits \(3\), and spoiling the top endpoint admits \(2\).
\end{example}

\noindent The congruences are the
\lword{Erdos1049/ZudilinConeArithmetic.lean}{203}{homEvalThreeTwo_mod_three}{bottom-endpoint
identity} and
\lword{Erdos1049/ZudilinConeArithmetic.lean}{320}{homEvalThreeTwo_mod_two}{top-endpoint
identity}; the unit consequences are the
\lword{Erdos1049/ZudilinConeArithmetic.lean}{338}{three_not_dvd_homEvalThreeTwo_of_const_unit}{constant-endpoint
obstruction} and
\lword{Erdos1049/ZudilinConeArithmetic.lean}{356}{two_not_dvd_homEvalThreeTwo_of_top_unit}{top-endpoint
obstruction}.

The proof is two lines, but the shape of the statement is not incidental.  The
difficulty lies in the fact that the specialisation sees a different endpoint at
each of the two primes: modulo \(3\) only the constant coefficient survives, and
modulo \(2\) only the top one does.  The two exclusions are therefore conditions
at opposite ends of the coefficient vector, and the statement below imposes one
at each end, on the two entries of a single coefficient pair.

\begin{proposition}[common divisor]\label{long1049:res:commonmult}
Let \(U,V\in\Z[X]\) and let \(W\ge0\).  If \(U\) has unit top endpoint, \(V\)
has unit constant endpoint, and an integer
\(c\) divides both \(H_W(U)\) and \(H_W(V)\), then
\[
 2\nmid c\qquad\text{and}\qquad 3\nmid c.
\]
\end{proposition}

\noindent This is the checked
\lword{Erdos1049/ZudilinConeArithmetic.lean}{398}{commonMultiplier_not_two_not_three_of_endpoint_units}{common-divisor
exclusion}.  The proposition does not say that the two evaluations are coprime;
it says that whatever they share misses both of the primes that matter at
\(3/2\).

\begin{example}\label{long1049:ex:commonmult}
Take \(W=2\), \(U=X^2+3\) and \(V=5X^2+1\).  The top endpoint of \(U\) and the
constant endpoint of \(V\) are both \(1\), and
\[
 H_2(U)=3\cdot4+1\cdot9=21,\qquad H_2(V)=1\cdot4+5\cdot9=49 .
\]
Here \(\gcd(21,49)=7\), so a common divisor does exist and is not small; it
is simply coprime to \(6\).
\end{example}

\begin{corollary}[no gain from integer scalar content or the stated common divisor at
\(3/2\)]\label{long1049:res:nomult}
Under the endpoint hypotheses of Proposition~\ref{long1049:res:commonmult}, neither
integer scalar content of specialised rows nor a common divisor of the two specialised
evaluations \(H_W(U)\) and \(H_W(V)\) can supply factors \(2\) and \(3\) by
those mechanisms in the common-width endpoint architecture studied here.
\end{corollary}

\begin{proof}
By Theorem~\ref{long1049:res:content} integer scalar content multiplies the analytic error and
the exterior determinant, including the absolute determinant height, by exactly
the factors it introduces, so a divisor obtained that way is paid for by the
same factor in that height.  By Proposition~\ref{long1049:res:commonmult} an integer
dividing both specialised evaluations is divisible by neither \(2\) nor \(3\).
\end{proof}

One further consequence of Theorem~\ref{long1049:res:endpoints} is worth stating,
because it bears on the most natural way one might hope to import an existing
denominator reduction.  Write $\Phi_m$ for the $m$th cyclotomic polynomial and,
for coprime $a>b\ge1$, put $\Phi_m(a,b)=b^{\varphi(m)}\Phi_m(a/b)$ for its
homogenisation at the declared width $\varphi(m)=\deg\Phi_m$.

\begin{proposition}[homogenised cyclotomic values are unit at both endpoints]
\label{long1049:res:cyclounit}
Let $a>b\ge1$ with $\gcd(a,b)=1$ and let $m\ge1$.  Then
$\gcd(\Phi_m(a,b),ab)=1$.  In particular $\gcd(\Phi_m(3,2),6)=1$ for every $m$.
\end{proposition}

\begin{proof}
$\Phi_m$ is monic and $\Phi_m(0)=\pm1$, so its coefficients at both declared
endpoints are units.  Theorem~\ref{long1049:res:endpoints} applies verbatim with $W$
replaced by $\varphi(m)$: modulo a prime dividing $b$ only the top term of the
homogenisation survives, and modulo a prime dividing $a$ only the constant term
does.
\end{proof}

The kernel-checked declaration
\lref{Erdos1049/ZudilinConeArithmetic.lean}{302}{cyclotomicHomEval_isCoprime_mul}
proves Proposition~\ref{long1049:res:cyclounit} in the same homogeneous-evaluation
representation.  It checks Proposition~3.6 under the displayed coprimality
assumptions; it does not certify the later analytic deductions or
Proposition~8.6.

The methods that reduce denominators at integer bases, the factorial-coset
quotients of Rhin and Viola~\cite{rhinviola1996} and the order-twelve group and
cyclotomic divisor of Zudilin~\cite{zudilin2004}, produce their gain as
cyclotomic or factorial factors of the coefficient polynomials.
Proposition~\ref{long1049:res:cyclounit} says that transporting such a factor through the
homogenisation at $(3,2)$ contributes no power of $2$ and no power of $3$,
whatever its size.  This does not make such factors useless: a large odd
divisor still reduces Archimedean height, and that is a different account of the
same product formula.  It does say that the $2$- and $3$-primary gain the
architecture of Section~\ref{long1049:sec:open} requires must come from somewhere other
than an imported cyclotomic factor.  The exclusions therefore reflect the
architecture itself, beyond the local behaviour of the two primes involved.

\noindent Both ingredients are Lean-checked; the combination is an ordinary
deduction and is not separately formalised.  The corollary excludes two ways of
producing the targeted gain.  It does not show that a gain of that kind is
necessary for a proof by linear forms at \(3/2\).

Corollary~\ref{long1049:res:nomult} excludes the two multiplicative mechanisms above,
and one candidate pursued in the rest of this section is additive: take an
integer combination of several rows and ask that the combination be divisible
where the individual rows are not.  No single row is multiplied by a scalar.  The endpoint
congruences are the first case of a divisibility condition that can be imposed
to any depth, and it is that condition, read additively, which is counted
below.

We first raise the two congruences to prime powers.  Fix depths \(R,S\ge0\).
For \(P\in\Z[X]\) the \emph{bottom jet} \(J_{3,R}(P)\) is the residue of
\(H_W(P)\) modulo \(3^R\), and the \emph{top jet} \(J_{2,S}(P)\) is its residue
modulo \(2^S\); Theorem~\ref{long1049:res:endpoints} computes them at \(R=S=1\).  Their
vanishing is exactly the requested divisibility:
\[
 J_{3,R}(P)=0\iff 3^R\mid H_W(P),\qquad
 J_{2,S}(P)=0\iff 2^S\mid H_W(P).
\]
These are the checked
\lword{Erdos1049/ZudilinConeArithmetic.lean}{191}{bottomJet3_eq_zero_iff_dvd}{bottom-jet
divisibility criterion} and
\lword{Erdos1049/ZudilinConeArithmetic.lean}{197}{topJet2_eq_zero_iff_dvd}{top-jet
divisibility criterion}.

The \emph{four-jet signature} of a coefficient pair
\((U,V)\) is then the quadruple
\[
 \bigl(J_{3,R}(U),J_{3,R}(V),J_{2,S}(U),J_{2,S}(V)\bigr)
 \in(\Z/3^R\Z)^2\times(\Z/2^S\Z)^2 ,
\]
two residues for each of the two primes, one from each entry of the pair.  By
the displayed criteria it vanishes exactly when \(3^R\) divides both specialised
entries and \(2^S\) divides both.

For this candidate architecture, this turns the targeted local divisor into an additive
congruence-kernel problem: what is sought is no longer a common divisor of
the two evaluations, but a vector of small integer coefficients on which four
residues vanish at once.  Since \(H_W\) is linear in the coefficients of \(P\),
the four-jet signature of a combination is the corresponding combination of
signatures, which is what makes the following count possible.

\begin{theorem}[binary four-jet collision]\label{long1049:res:jetkernel}
Fix a width \(W\) and depths \(R,S\), and let \((U_j,V_j)_{j<M}\) be any \(M\)
pairs of integral polynomials.  Call a subset of \(\{0,\dots,M-1\}\), equivalently
a vector of \(\{0,1\}^M\), a \emph{binary selector}.  If the
\(2^M\) binary selectors outnumber the finite four-jet target
\[
 (\Z/3^R\Z)^2\times(\Z/2^S\Z)^2,
\]
then two distinct subsets have the same four-jet sum.  Subtracting their
indicator vectors gives a nonzero coefficient vector in
\(\{-1,0,1\}^M\) cancelling all four jets.  The target has exact cardinality
\[
 (3^R)^2(2^S)^2.
\]
In particular, if \(R>0\) and \(4R+2S\le M\), such a collision exists.
\end{theorem}

\begin{proof}
Send each binary selector to the sum of the four-jet signatures it selects.
The claimed cardinal inequality and the pigeonhole principle give two
distinct selectors in the same fibre.  The cardinality formula is the product
of the four cyclic-modulus cardinalities.  For \(R>0\),
\[
 (3^R)^2(2^S)^2<(4^R)^2(2^S)^2=2^{4R+2S}\le2^M,
\]
which proves the stated sufficient threshold.
\end{proof}

The power bracket improves the generic coefficient $4R$ when the bottom depth
is a multiple of $41$.

\begin{corollary}[sharp rank-$41$ four-jet threshold]\label{long1049:res:rankfortyone}
Let $T>0$.  At bottom depth $R=41T$, any family of
$M\ge130T+2S$ integral polynomial pairs has two distinct binary selectors
with the same four-jet sum.  For $T=1$ the coefficient $130$ is exact for
this counting argument:
\[
 2^{129+2S}<\bigl| (\Z/3^{41}\Z)^2\times
                       (\Z/2^S\Z)^2\bigr|.
\]
No exact-optimality assertion is made here for $T>1$.
\end{corollary}

\begin{proof}
The upper power inequality gives
\[
 (3^{41T})^2(2^S)^2
 <(2^{65})^{2T}(2^S)^2=2^{130T+2S}\le2^M,
\]
so Theorem~\ref{long1049:res:jetkernel} applies.  For $T=1$, direct integer
evaluation gives $2^{129}<3^{82}$; multiplying by $(2^S)^2$ gives the
displayed reverse count at $129+2S$.
\end{proof}

\noindent The sufficient collision threshold is
\lword{Erdos1049/AdelicHeightBridge.lean}{1800}{exists_distinct_binary_selectors_same_fourJet_of_rank_41}{checked here},
while its exact unit-block failure one row earlier is
\lword{Erdos1049/AdelicHeightBridge.lean}{1813}{fourJet_card_gt_two_pow_of_rank_41}{checked here}.

For all depths at once the ambient-cardinality inequality $2^{M}>3^{2R}2^{2S}$
holds exactly when
\[
 M\ \ge\ M_{\mathrm{count}}(R,S)=\bigl\lfloor2R\log_23+2S\bigr\rfloor+1,
\]
because $2R\log_23$ is irrational for $R\ge1$.  At $R=41$ this is $130+2S$, and
at $R=41\cdot31$ it is $4029+2S$, one below the uniform bound $4030+2S$ of
Corollary~\ref{long1049:res:rankfortyone}; the corollary trades exactness for a
certificate that is a single integer comparison.  The name of the parameter $41$
is the bottom depth, and the number of input rows is $130+2S$.

Counting alone does not ensure that the two selectors produce different
analytic remainders.  The exact missing step is a bounded-fibre estimate.

\begin{theorem}[bounded-fibre escape]\label{long1049:res:boundedfibre}
Let $A$ and $B$ be finite sets, let $f:A\to B$, and let $g:A\to C$ be any
map into a set $C$.  Suppose every fibre of $g$ has at most $k$ elements.  If
\[
 |B|k<|A|,
\]
then there exist distinct $x,y\in A$ such that
\[
 f(x)=f(y)\qquad\hbox{and}\qquad g(x)\ne g(y).
\]
Thus, with $f$ the four-jet sum and $g$ the analytic remainder, a uniform
remainder-multiplicity bound converts surplus selector entropy into a
four-jet collision outside the remainder nullspace.
\end{theorem}

\begin{proof}
If every pair in a common $f$-fibre also had the same $g$-value, each
$f$-fibre would lie in one $g$-fibre and hence have size at most $k$.
Summing over the at most $|B|$ fibres of $f$ would give
$|A|\le |B|k$, contrary to the hypothesis.
\end{proof}

\noindent This finite escape principle is
\lword{Erdos1049/AdelicHeightBridge.lean}{1832}{exists_ne_map_eq_map_ne_of_card_mul_lt}{checked here}.

\begin{theorem}[B\'ezout--Pl\"ucker tail collapse]
\label{long1049:res:plucker-collapse}
Let $R_0$ be a commutative ring and let $w_n=(A_n,B_n)\in R_0^2$.  Suppose
that each row is unimodular, meaning that $u_nA_n+v_nB_n=1$ for some
$u_n,v_n\in R_0$, and that every adjacent minor vanishes:
\[
 A_nB_{n+1}-B_nA_{n+1}=0\qquad(n\ge0).
\]
Then every pairwise minor $A_iB_j-B_iA_j$ vanishes.  In particular, take
$R_0=\Z/(2^S3^R)\Z$ with $R>0$.  If $S+2R\le k$, there are two distinct
binary selectors $s,t\in\{0,1\}^{k}$ such that
\[
 \sum_{i<k}s_iw_i=\sum_{i<k}t_iw_i.
\]
Thus the sufficient width is $S+2R$.  The ambient two-coordinate width is
$2S+4R$.
\end{theorem}

\begin{proof}
If $(a,b)$ is unimodular, say $ua+vb=1$, and $ay-bx=0$, then
\[
 (x,y)=(ux+vy)(a,b).
\]
Indeed, the first coordinate follows by replacing $ay$ with $bx$, and the
second by the reverse substitution.  Apply this identity to consecutive rows:
each next row is a scalar multiple of the current one.  Induction places the
entire tail on the line through $w_0$ and proves the pairwise-minor assertion.
A determinant-one B\'ezout shear sends $w_0$ to
$(1,0)$, so all selector sums have only one free residue coordinate and occupy
at most $2^S3^R$ values.  Finally
\[
 2^S3^R<2^S4^R=2^{S+2R}\le2^k,
\]
and pigeonhole gives the two selectors.
\end{proof}

\noindent No particular coordinate needs to be invertible: modulo six,
$(2,3)$ is unimodular because $-2+3=1$, although neither entry is a unit.
Some nondegeneracy is essential.  The three rows $(1,0),(0,0),(0,1)$ have
zero adjacent minors but outer minor one; a zero middle row transmits no
information between its neighbours.

\noindent The unimodular
\href{https://github.com/wcook04/plectis-erdos/blob/0cfa24a7fe555d75a9d9e7f119da4720a88c1396/ErdosProblems/Erdos1049/BezoutPluckerJets.lean\#L191}{adjacent-to-pairwise determinant propagation}
and the resulting
\href{https://github.com/wcook04/plectis-erdos/blob/0cfa24a7fe555d75a9d9e7f119da4720a88c1396/ErdosProblems/Erdos1049/BezoutPluckerJets.lean\#L280}{modular selector collision}
are Lean-checked, including the
\href{https://github.com/wcook04/plectis-erdos/blob/0cfa24a7fe555d75a9d9e7f119da4720a88c1396/ErdosProblems/Erdos1049/BezoutPluckerJets.lean\#L299}{explicit $S+2R$ threshold}.
The unit-coordinate statements are special cases of the stronger
rowwise-coprime theorem.
This conditional theorem is stronger than the ambient four-jet count only
after its minor-vanishing hypothesis has been established.  No such all-tail
hypothesis is proved here for an actual $q$-Ap\'ery or Zudilin family, and the
theorem says nothing about whether the resulting selector difference has
nonzero analytic remainder.

\noindent\textbf{Boundary.}
Corollary~\ref{long1049:res:rankfortyone} is a sharp finite kernel statement at
$T=1$, not an analytic nonvanishing theorem.  Theorem~\ref{long1049:res:boundedfibre}
identifies the precise extra input needed to escape the nullspace, but this
paper does not prove a multiplicity bound for the actual $q$-Ap\'ery or
Zudilin remainder family.

\noindent The target count is the checked
\lword{Erdos1049/ZudilinConeArithmetic.lean}{131}{fourJetSignature_card}{four-jet
target cardinality}; the abstract collision is the checked
\lword{Erdos1049/ZudilinConeArithmetic.lean}{142}{exists_distinct_binary_selectors_same_fourJet}{four-jet
pigeonhole kernel}, and the linear sufficient condition is the checked
\lword{Erdos1049/ZudilinConeArithmetic.lean}{161}{exists_distinct_binary_selectors_same_fourJet_of_rank}{rank--depth
collision threshold}.  The count is a routine pigeonhole; the reformulation
above is what makes it relevant.  Pigeonhole cancellation itself requires no
independence.  Additional information about the input family is needed to
ensure that the resulting nonzero selector difference has a nonzero combined
polynomial pair and analytic remainder.  None of the statements proved here
supplies such a family or proves either nonvanishing conclusion.

\begin{example}\label{long1049:ex:jetcount}
At depths \(R=S=1\) the four-jet target is
\((\Z/3\Z)^2\times(\Z/2\Z)^2\), of cardinality \(9\cdot4=36\), and the threshold
reads \(M\ge4\cdot1+2\cdot1=6\).  With six pairs there are \(2^{6}=64\) binary
selectors against \(36\) targets, so two of them collide and their difference
is a vector in \(\{-1,0,1\}^{6}\), not identically zero, killing all four jets.
\end{example}

Informally, the theorem says only this: once there are at least \(4R+2S\) pairs
and the bottom depth \(R\) is positive, some coefficient vector in
\(\{-1,0,1\}^M\), not identically zero, kills all four jets of the
corresponding combination.  It does not say that the combination is nonzero as
a pair of polynomials, and it does not say that its remainder is nonzero.  Those
are the two obligations Problem~\ref{long1049:prob:kernel} carries.

One further exclusion is recorded here.  Its subject is the scalar parameters
of Zudilin's cone, and it says nothing about the coefficient polynomials.

\begin{theorem}[scalar margin no-go]\label{long1049:res:scalar}
Let \(C_1>0\).  If \(C_0\le0\) or \(2C_0\le C_1\), then
\[
 C_0\log 3-C_1\log 2<0.
\]
In the positive branch $C_0>0$ and $2C_0\le C_1$, the stronger estimate is
\[
 C_0\log3-C_1\log2< -\frac{17}{41}C_0\log2.
\]
\end{theorem}

\noindent Written multiplicatively, the conclusion is \(3^{C_0}<2^{C_1}\).  The
inequality is immediate from \(\log3<2\log2\) and \(C_1>0\), and
is the checked
\lword{Erdos1049/ZudilinConeArithmetic.lean}{428}{three_two_scalar_margin_neg}{three-halves
scalar margin}.  The positive-branch deficit is the checked
\lword{Erdos1049/AdelicHeightBridge.lean}{1763}{three_two_scalar_margin_lt_explicit}{$17/41$ margin}.
The interest is in the fence around it.  The primary Zudilin
theorem~\cite[Theorem~1, p.~154; Sec.~5, pp.~161--162]{zudilin2004} supplies an integer-base
irrationality-exponent estimate on its parameter cone, and the elementary
inequality \(\mu\ge2\) then forces \(2C_0\le C_1\) whenever \(C_0>0\); Lean
checks that implication separately.  The primary theorem assumes an integer
base \(p=1/q\).  It does not state a rational \(p=3/2\) theorem, so the
all-scale coefficient construction and rational specialisation remain
external to the checked result, which is the scalar parameter margin alone.

\section{Failure of coordinatewise clearing at $3/2$}\label{long1049:sec:corridor}

This section and the next return to the elementary route and record what the
residue $s^{n}$ costs there.  We first isolate the arithmetic that the clearing
scheme leaves behind, in a form that does not mention the series.  The point of
doing so is that the clearing argument asks for two things at once, that a
power of the numerator divide what has been accumulated and that what survives
be small, and these are easier to play off against each other once the series
has been discarded and only six natural numbers remain.

\begin{definition}[coordinatewise corridor]\label{long1049:def:corridor}
Let $a,b,N,K,Q,D$ be natural numbers.  Say that $(a,b,N,K,Q,D)$ is a
\emph{coordinatewise corridor} when
\[
 a>0,\qquad Q>0,\qquad D>0,\qquad D\le N+K,\qquad
 a^{K}\mid QD,\qquad Q\,b^{\,N+K+1}<a^{\,K+1}.
\]
\end{definition}

\noindent The reading is: $a$ and $b$ are the numerator and denominator of the
base, playing the roles of $r$ and $s$ in the introduction, so that
$(a,b)=(3,2)$ is the case of interest;
$N$ is the shift, $K$ is the width of the cleared window, $Q$ is the
accumulated clearing factor, and $D$ is the final coefficient being cleared.
The bound $D\le N+K$ is the only property of the coefficient used; for the
divisor-counting coefficient it holds because $\tau(n)\le n$.  The divisibility
$a^{K}\mid QD$ is the requirement that clearing succeeded coordinatewise, and
the last inequality is the tail estimate that makes the trapped integer smaller
than~$1$.  The name records the shape of the constraint: the divisibility
bounds $a^{K}$ from above by $Q(N+K)$, the tail estimate bounds $a^{K+1}$ from
below by $Q\,b^{\,N+K+1}$, and admissible parameters must fit in the band
between them.  The definition is the
\lword{Erdos1049/RationalBaseLambert.lean}{149}{CoordinatewiseCorridor}{corridor
predicate}.

\begin{theorem}[corridor bound]\label{long1049:res:corridorbound}
If $(a,b,N,K,Q,D)$ is a coordinatewise corridor, then
\[
 b^{\,N+K+1}<a\,(N+K).
\]
\end{theorem}

\begin{proof}
A routine divisibility computation.  Since $QD>0$ and $a^{K}\mid QD$, we have
$a^{K}\le QD$, and $D\le N+K$ gives
$a^{K}\le Q(N+K)$.  Hence
\[
 Q\,b^{\,N+K+1}<a^{\,K+1}=a^{K}\cdot a\le Q(N+K)\cdot a,
\]
and cancelling the positive factor $Q$ gives the claim.
\end{proof}

\noindent Formalised as the
\lword{Erdos1049/RationalBaseLambert.lean}{157}{coordinatewiseCorridor_implies_pow_lt_linear}{power-versus-linear
consequence}.

The inequality of Theorem~\ref{long1049:res:corridorbound} is where the integer and
rational cases part.  At $b=1$ it reads $1<a(N+K)$, which holds for every
$a\ge2$ and every nonempty window; this necessary inequality imposes no
obstruction.  The other corridor hypotheses remain in force.
At $b\ge2$ the left side is exponential in $N+K$ and the right side is linear,
so the corridor can survive only for small $N+K$.  At $(a,b)=(3,2)$ the
crossing has already happened at the smallest admissible window.

\begin{example}\label{long1049:ex:corridor}
Take the smallest window, $N=K=1$, and numerator $a=3$.  At $b=1$ the tuple
$(3,1,1,1,3,1)$ is a corridor: $D=1\le2$, the divisibility reads $3\mid3$, and
the tail estimate reads $3\cdot1^{3}=3<3^{2}=9$.  At $b=2$ no choice works.  The
tail estimate becomes $Q\cdot2^{3}<3^{2}$, which forces $Q=1$; the divisibility
then reads $3\mid D$, and the only candidates are $D=1$ and $D=2$, neither
divisible by $3$.  This is the proof of
Theorem~\ref{long1049:res:corridorbound} in miniature: it derives $a^{K}\le Q(N+K)$,
here $3\le2$.
\end{example}

\begin{proposition}\label{long1049:res:exp}
For every natural number $x\ge2$ we have $3x<2^{\,x+1}$.
\end{proposition}

\begin{proof}
A routine induction from $x=2$, where $6<8$.  For the step, $2^{x+1}\ge2^{2}>3$ when
$x\ge1$, so $3(x+1)=3x+3<2^{x+1}+2^{x+1}=2^{x+2}$.
\end{proof}

\noindent Formalised as the
\lword{Erdos1049/RationalBaseLambert.lean}{178}{three_mul_lt_two_pow_succ}{exponential
comparison}.

\begin{theorem}[no corridor at base $3/2$]\label{long1049:res:nocorridor}
For all $N\ge1$ and $K\ge1$ and all natural $Q,D$, the tuple
$(3,2,N,K,Q,D)$ is not a coordinatewise corridor.
\end{theorem}

\begin{proof}
A corridor would give $2^{\,N+K+1}<3(N+K)$ by
Theorem~\ref{long1049:res:corridorbound}, contradicting Proposition~\ref{long1049:res:exp}
applied to $x=N+K\ge2$.
\end{proof}

\noindent Formalised as the
\lword{Erdos1049/RationalBaseLambert.lean}{191}{threeHalves_no_coordinatewiseCorridor}{corridor
exclusion at three halves}.

Theorem~\ref{long1049:res:nocorridor} excludes the coordinatewise clearing scheme at
$3/2$, and nothing else: it does not bound the denominator of $F(3/2)$, it does
not show that $F(3/2)$ is irrational, and it does not show that $F(3/2)$ is
rational.  It also does not cover a clearing scheme of a different shape, since
the corridor fixes one divisibility pattern and one tail inequality.

\section{The cleared-tail recurrence and the size of the forcing
term}\label{long1049:sec:tail}

Theorem~\ref{long1049:res:nocorridor} says that one scheme fails.  This section
identifies the quantity responsible, as an exact recurrence.

Let $r,s,B,F$ be rationals with $r\ne0$ and let $c:\N\to\Q$ be arbitrary.
Define the prefix and the cleared tail state by
\[
 P_N=\sum_{m=0}^{N-1}\frac{c(m+1)\,s^{\,m+1}}{r^{\,m+1}},
 \qquad
 U_N=B\,r^{N}\bigl(F-P_N\bigr).
\tag{$\ast$}\label{long1049:eq:tailstate}
\]
Thus $P_N$ is the partial sum of $\sum_{n\ge1}c(n)(s/r)^{n}$ through level $N$,
and $U_N$ is the tail of a putative value $F$ after that level, scaled by
$Br^{N}$.  These are the
\lword{Erdos1049/RationalBaseLambert.lean}{204}{rationalBasePrefixQ}{rational-base
prefix} and the
\lword{Erdos1049/RationalBaseLambert.lean}{217}{rationalBaseClearedTailQ}{cleared
tail state}.

\begin{theorem}[cleared-tail recurrence]\label{long1049:res:tailrec}
Let $r,s,B,F\in\Q$ with $r\ne0$, let $c:\N\to\Q$, and let $P_N$ and $U_N$ be as
in~\eqref{long1049:eq:tailstate}.  Then for every $N$,
\[
 U_{N+1}=r\,U_N-B\,c(N+1)\,s^{\,N+1}.
\]
\end{theorem}

\begin{proof}
An immediate computation.  Expanding
$P_{N+1}=P_N+c(N+1)s^{N+1}/r^{N+1}$ and $r^{N+1}=r^{N}\cdot r$ in the
definition of $U_{N+1}$ and clearing the denominator $r^{N+1}$, which is
nonzero, gives the identity.
\end{proof}

\noindent Formalised as the
\lword{Erdos1049/RationalBaseLambert.lean}{223}{rationalBaseClearedTailQ_succ}{cleared-tail
recurrence}.

The recurrence has a linear part $rU_N$ and a forcing term $Bc(N+1)s^{N+1}$,
the inhomogeneous term the state receives at each step.
The direct-clearing route studied here works only if the state remains in a
bounded window, and the forcing term is what that window must absorb.  Its
size is what separates the two denominator regimes.

\begin{theorem}[the forcing term]\label{long1049:res:forcing}
Let $s,B$ be natural numbers and $c:\N\to\N$, and put
$G_N=B\,c(N+1)\,s^{\,N+1}$.
\begin{enumerate}
\item If $s\ge2$, $B\ge1$ and $c(N+1)\ge1$, then $2^{\,N+1}\le G_N$.
\item If $s=1$, then $G_N=B\,c(N+1)$.
\end{enumerate}
\end{theorem}

\begin{proof}
Both parts are routine.  For the first,
$2^{N+1}\le s^{N+1}=1\cdot s^{N+1}\le Bc(N+1)s^{N+1}$,
using $B\,c(N+1)\ge1$.  The second part is the definition with $s=1$.
\end{proof}

\noindent Formalised as the
\lword{Erdos1049/RationalBaseLambert.lean}{240}{twoPow_le_rationalBaseForcingNat}{exponential
lower bound} and the
\lword{Erdos1049/RationalBaseLambert.lean}{254}{rationalBaseForcingNat_one}{integer-base
collapse}.

At $s=1$ the forcing term is $Bc(N+1)$, so it grows only as fast as the
coefficient; for the divisor function that is $O(N^{\varepsilon})$ for every
$\varepsilon>0$, and a bounded-state argument has room.  At $s\ge2$ the same
term is at least $2^{N+1}$ whenever the coefficient is nonzero.

\begin{example}\label{long1049:ex:forcing}
Take $B=1$, $c=\tau$ and $N=9$, so that the coefficient is $\tau(10)=4$.  At
$s=1$ the forcing term is $4$.  At $s=2$ it is $4\cdot2^{10}=4096$, and part~(1)
of Theorem~\ref{long1049:res:forcing} already guarantees at least $2^{10}=1024$ without
knowing the coefficient at all.
\end{example}

This is an
exact lower bound on one quantity, and it is all that is proved.  It is not a
proof that no bounded-state argument exists at $s\ge2$; the theorem that one
particular scheme fails is Theorem~\ref{long1049:res:nocorridor}, and the bound here
records the size of the term that scheme would have to absorb.

\section{The height criterion at $7/2$}\label{long1049:sec:sevenhalves}

In 1994 Bundschuh and V\"a\"an\"anen proved an irrationality criterion for a
family of rational bases cut out by a height condition
\cite[Theorem~2, p.~177; hypotheses pp.~175--176]{bv1994}.  We keep their
notation: $q$ is the base, and $\alpha$
and $\lambda$ are the parameters of the criterion.  In its special case
$\alpha=-1$, the printed hypothesis is
\[
 \lambda<\left(\frac12+\frac1{\pi^2}\right)^{-1}.
\]
At $q=7/2$ the Archimedean parameter is
$\lambda=\log 7/\log(7/2)$.  The criterion therefore applies once the following
strict inequality is checked.

\begin{theorem}[the $7/2$ height condition]\label{long1049:res:sevenhalves}
\[
 \frac{\log 7}{\log(7/2)}
 <
 \left(\frac12+\frac1{\pi^2}\right)^{-1}.
\]
\end{theorem}

\noindent Numerically the two sides are $1.5533\ldots$ and $1.6630\ldots$, so
the condition holds with a margin of about $0.11$.  The Lean proof factors the
estimate through the explicit
\lword{Erdos1049/RationalBaseLambert.lean}{44}{BundschuhVaananenHeightRegion}{height-region
predicate} and the
\lword{Erdos1049/RationalBaseLambert.lean}{32}{sevenHalves_power_certificate}{integer
certificate} $2^{18}<7^7$, the resulting
\lword{Erdos1049/RationalBaseLambert.lean}{48}{sevenHalves_log_ratio_lt}{logarithmic
ratio bound} $\log2/\log7<7/18$, the
\lword{Erdos1049/RationalBaseLambert.lean}{64}{one_div_pi_sq_lt_one_nine}{$\pi$-bound}
$1/\pi^2<1/9$, and the
\lword{Erdos1049/RationalBaseLambert.lean}{73}{sevenHalves_bundschuhVaananen_margin}{strict
margin}
\[
 \frac{\log2}{\log7}<\frac12-\frac1{\pi^2}.
\]
The
\lword{Erdos1049/RationalBaseLambert.lean}{119}{sevenHalves_archimedean_height_condition}{final
height inequality} then rewrites
$\log(7/2)=\log7-\log2$ and closes the displayed condition.

This formalises the complete elementary parameter check at $q=7/2$; it does
not formalise Bundschuh and V\"a\"an\"anen's analytic irrationality theorem,
whose proof occupies pp.~189--193 of the source.  The conclusion that
$F(7/2)$ is irrational is consequently cited from that theorem, not claimed as
a Lean theorem here.

The criterion does not cover $3/2$, and the results of
Sections~\ref{long1049:sec:corridor} and~\ref{long1049:sec:tail} give no evidence either way
about whether $F(3/2)$ is irrational.

\subsection{The $81/200$ logarithmic region}

Consider the rational-height region
\[
 \frac{\log b}{\log a}<\frac{81}{200}
 \qquad(a>b>0).
\]
The rational number $81/200$ is not itself a threshold with an analytic theorem
attached to it.  It is the lower half of the rational bracket
$81/200<\theta^{*}<1/2$ of Section~\ref{long1049:sec:regionbracket} around the derived
constant of Theorem~\ref{long1049:res:region}, and its role is to make membership of the
region of that theorem decidable by an integer comparison.  The Lean module of
this library treats the displayed inequality as a definition and proves
elementary memberships and exclusions.  In particular,
\[
 \frac25
 <\frac{\log4}{\log31}
 <\frac{81}{200}
 <\frac{\log2}{\log3}.
\]
The first two comparisons come respectively from $31^2<4^5$ and
$4^{200}<31^{81}$; the last comes from $3^{81}<2^{200}$.  The lower bound is
the checked theorem
\lword{Erdos1049/ZudilinHeightRegion.lean}{65}{twoFifths_lt_thirtyoneFour_log_ratio}{two-fifths
lower bound}.  Consequently $31/4$, and every positive power $(31/4)^r$, lies
in the enlarged region
(\lword{Erdos1049/ZudilinHeightRegion.lean}{45}{thirtyoneFour_mem_zudilinHeightRegion}{checked},
\lword{Erdos1049/ZudilinHeightRegion.lean}{59}{thirtyoneFour_power_mem_zudilinHeightRegion}{power family}).
The same base lies strictly outside the earlier
Bundschuh--V\"a\"an\"anen region
(\lword{Erdos1049/ZudilinHeightRegion.lean}{91}{thirtyoneFour_outside_bundschuhVaananenHeightRegion}{checked}),
so the $81/200$ inequality defines a strict set-theoretic enlargement of the
earlier recorded logarithmic region.  Combined with the bracket
$81/200<\theta^{*}$, these memberships are exactly the finite input to
Theorem~\ref{long1049:res:31over4}, which is where the irrationality of $F(31/4)$ and of
every $F\bigl((31/4)^{r}\bigr)$ is proved.

It still does not approach $3/2$.  The exact comparison
\[
 3^{81}<2^{200}
 \quad\Longrightarrow\quad
 \frac{81}{200}<\frac{\log2}{\log3}
\]
is Lean-checked, as is the conclusion that $3/2$ belongs to neither height
region
(\lword{Erdos1049/ZudilinHeightRegion.lean}{105}{eightyOneTwoHundredths_lt_threeHalves_log_ratio}{boundary},
\lword{Erdos1049/ZudilinHeightRegion.lean}{122}{threeHalves_outside_zudilinHeightRegion}{exclusion}).
Thus $81/200$ is the larger of the two explicitly defined cutoffs used below,
while a cutoff that includes $3/2$ must be strictly larger than
$\log2/\log3\approx0.6309$.

\section{Denominator exponents for a homogenised Pad\'e
construction}\label{long1049:sec:pade}

A second external route builds the linear forms of Section~\ref{long1049:sec:primitive}
by Pad\'e approximation.  Such a construction produces its coefficients as sums
of rational summands, and to obtain integer rows one multiplies through by a
single common denominator.  The bookkeeping obligation is then to check that no
summand needs a larger denominator than the one proposed, which is a comparison
between the exponents alone.  This section does that comparison, and only that.

No source construction is named here for the exponent expressions below, and
they are checked as displayed polynomial identities.  No derivation of them is
given here.
Giving them Pad\'e-theoretic force would additionally require deriving the
expressions from a stated homogenised coefficient formula, proving integrality
of the coefficients after multiplication by the proposed common denominator,
and establishing nonvanishing and decay of the associated remainder.  None of
those three is claimed here.

For a
homogenised construction over integer parameters the proposed common
denominator exponent is $E_n=(3n^{2}-n)/2$.  Since only doubled exponents
occur below, every statement lives over $\Z$ and no parity bookkeeping is
needed; we write $\widetilde{E}_n=2E_n=3n^{2}-n$.

\begin{proposition}[summand bound and exact gap]\label{long1049:res:pade}
Let $\widetilde{E}_n=3n^{2}-n$ and put
\[
 \widetilde{P}(n,k)=2\bigl(k(n-k)+nk\bigr)+k(k-1),
\]
\[
 \widetilde{Q}(n,m)=2(n^{2}-n)+j^{2}+2jm+j-m^{2}+3m,
 \qquad j=n-m-1 .
\]
Then, for integers $n,k,m$:
\begin{enumerate}
\item if $0\le k\le n$, then $\widetilde{P}(n,k)\le\widetilde{E}_n$, and the
gap factors as
$\widetilde{E}_n-\widetilde{P}(n,k)=(n-k)(3n-k-1)$;
\item $\widetilde{E}_n-\widetilde{Q}(n,m)=2\bigl(n+m(m-1)\bigr)$ identically;
\item if $n\ge0$ and $m\ge1$, then $\widetilde{Q}(n,m)\le\widetilde{E}_n$.
\end{enumerate}
\end{proposition}

\noindent Part~(1) is the
\lword{Erdos1049/RationalPadeArithmetic.lean}{30}{rationalPadePSummandDenExpTwice_le}{summand
exponent bound}, part~(2) the
\lword{Erdos1049/RationalPadeArithmetic.lean}{52}{rationalPadeQMaxDenExpTwice_gap}{exact
gap identity}, and part~(3) the
\lword{Erdos1049/RationalPadeArithmetic.lean}{60}{rationalPadeQMaxDenExpTwice_le}{maximal
exponent bound}.  The gap in~(2) is an identity in $\Z[n,m]$, so~(3) follows
from $m(m-1)\ge0$ and $n\ge0$; the hypothesis $m\ge1$ names the range in which
the corresponding summand does not vanish, and it is not the sharpest
hypothesis under which the inequality holds.

\begin{example}\label{long1049:ex:pade}
At $n=2$ the proposed doubled exponent is $\widetilde{E}_2=10$, and
$\widetilde{P}(2,k)$ takes the values $0,6,10$ at $k=0,1,2$.  The three gaps are
$10,4,0$, matching the factorisation $(2-k)(5-k)$ of part~(1); the summand at
$k=2$ is the one that saturates the proposed denominator.  For part~(2), at
$m=1$ we have $j=0$ and $\widetilde{Q}(2,1)=6$, with gap
$4=2\bigl(2+1\cdot0\bigr)$.
\end{example}

These are routine inequalities between polynomials in the exponents.  They
establish that the proposed exponent $\widetilde{E}_n$ dominates the two
displayed summand exponent expressions, and nothing further.  Positivity of the remainder, its rate of
decay, and the comparison of that rate against the denominator height are the
analytic obligations, and none of them is treated here, so nothing in this
section is an irrationality measure.

\section{Complements and further questions}\label{long1049:sec:open}

Problem~\#1049 remains open.  Theorem~\ref{long1049:res:jetkernel} suggests the following precise sufficient
subproblem for one common-width additive architecture.  It is not asserted to
be necessary for every proof of irrationality.

\begin{problem}[common-width simultaneous endpoint-jet construction]\label{long1049:prob:kernel}
Exhibit an integer constant \(C\ge1\) and, for every sufficiently large
positive integer \(n\), positive integers \(W_n,R_n,S_n,M_n\) such that
\[
 n^2\le W_n,R_n,S_n\le Cn^2,
 \qquad 4R_n+2S_n\le M_n\le Cn^2,
\]
together with polynomial pairs
\((U_{n,j},V_{n,j})\in\Z[X]^2\) for \(0\le j<M_n\), each of degree at most
the common declared width \(W_n\), whose specialised integer rows are
primitive:
\[
 \gcd\!\bigl(H_{W_n}(U_{n,j}),H_{W_n}(V_{n,j})\bigr)=1.
\]
Find a nonzero vector
\(\lambda^{(n)}\in\{-1,0,1\}^{M_n}\) for which, on putting
\[
 U_n=\sum_{j<M_n}\lambda^{(n)}_jU_{n,j},
 \qquad V_n=\sum_{j<M_n}\lambda^{(n)}_jV_{n,j},
\]
the pair \((U_n,V_n)\) is not \((0,0)\), all four common-width jets vanish,
\[
 J_{3,R_n}(U_n)=J_{3,R_n}(V_n)=0,
 \qquad J_{2,S_n}(U_n)=J_{2,S_n}(V_n)=0,
\]
where every jet in this display is formed using the declared width \(W_n\),
and the resulting divided integer linear form
\[
 A_n=\frac{H_{W_n}(U_n)}{3^{R_n}2^{S_n}},
 \qquad
 B_n=\frac{H_{W_n}(V_n)}{3^{R_n}2^{S_n}},
 \qquad
 \rho_n=A_nF(3/2)-B_n
\]
satisfies the explicit analytic condition
\[
 0<|\rho_n|<\frac1n.
\]
\end{problem}

The jet equations make \(A_n,B_n\) integers.  A solution would prove
irrationality: if \(F(3/2)=a/b\) in lowest terms, every nonzero \(\rho_n\) has
absolute value at least \(1/b\), contradicting the displayed bound for
\(n>b\).  Theorem~\ref{long1049:res:jetkernel} supplies only a nonzero signed relation
with the four jet equations once the pairs and size inequality are present; it
does not supply primitive input rows, a nonzero combined polynomial pair, or
the nonvanishing and decay of \(\rho_n\).

Integer rescaling and a common divisor of the two specialised evaluations are
the two mechanisms excluded by Corollary~\ref{long1049:res:nomult}.  Polynomial factors
before specialisation, cross-row or determinant-specific arithmetic,
cyclotomic factors, other additive constructions, and entirely different
architectures remain outside that corollary and are not decided either way.

The remaining linear-form argument therefore has two
independent gates: first find a four-jet collision that does not collapse
algebraically or analytically; only then ask whether its divisibility and decay
beat its height.  We state those gates separately below.

\subsection{Exact kernel escape}

Fix, for each $n$, a declared width $W_n$ and a specified family
\[
 (U_{n,j},V_{n,j},\mathcal R_{n,j})_{j<M_n},
\]
where $U_{n,j},V_{n,j}\in\Z[X]$ and $\mathcal R_{n,j}$ is the corresponding
remainder function.  Normalisation must be fixed before the jet map is formed.
Call the family \emph{polynomially primitive} when the common coefficient
content of the pair in $\Z[X]^2$ is divided out before specialisation; as the
Terminology paragraph records, that is a different operation from primitive
normalisation of a specialised integer row, which is what
Problem~\ref{long1049:prob:kernel} imposes.

The specialised row content
\[
 g^{\mathrm{spec}}_{n,j}
 =\gcd\!\bigl(H_{W_n}(U_{n,j}),H_{W_n}(V_{n,j})\bigr)
\]
and the content of a final additive combination are separate quantities; they
are not silently divided out.  If either is divided out later, both the height
and the remainder below are measured after that same division.  Under the unit
endpoint hypotheses, $g^{\mathrm{spec}}_{n,j}$ is coprime to $6$, so such
division does not manufacture or destroy the $2$- and $3$-primary jet
vanishing.

Dividing a specialised row by its content changes the polynomial family it came
from, and lifting the divided row back is a constrained problem.  At width $W$ the map $P\mapsto H_W(P)$ on integer polynomials of
degree at most $W$ is surjective onto $\Z$, since $3^{W}$ and $2^{W}$ are
coprime, so a lift always exists; what is not automatic is a lift carrying the
intended remainder identity and a uniform coefficient-height bound.  At width
$1$, for instance, $(X+1,X+1)$ specialises to $(5,5)$, whose primitive
normalisation $(1,1)$ lifts to $(X-1,X-1)$ and not to any rescaling of the
original pair.  Any lift used below is therefore supplied together with its
width, its height bound and its exact remainder.

For target depths $R_n,S_n$, let $J_n(\lambda)$ be the four-jet signature of
the pair $\sum_j\lambda_j(U_{n,j},V_{n,j})$, and define
\[
 \mathcal C_n=
 \{\lambda\in\{-1,0,1\}^{M_n}\mathbin{\backslash}\{0\}:J_n(\lambda)=0\},
\]
\[
\begin{aligned}
 K_n^{\mathrm{poly}}
   &=\left\{\lambda\in\{-1,0,1\}^{M_n}:
       \sum_j\lambda_jU_{n,j}=0,\ \sum_j\lambda_jV_{n,j}=0\right\},\\
 K_n^{\mathrm{rem}}
   &=\left\{\lambda\in\{-1,0,1\}^{M_n}:
       \sum_j\lambda_j\mathcal R_{n,j}(3/2)=0\right\}.
\end{aligned}
\]
The second nullspace concerns the value at $3/2$; an identically zero
remainder function is a still stronger collapse and is automatically bad.

The exact pigeonhole condition is
\[
 2^{M_n}>3^{2R_n}2^{2S_n},
 \qquad\text{equivalently}\qquad
 M_n>2R_n\log_2 3+2S_n.
\tag{8.1}\label{long1049:eq:exact-jet-threshold}
\]
Thus the least admissible integer rank is
\[
 M_{\min}(R,S)=
 \left\lfloor2R\log_2 3+2S\right\rfloor+1.
\]
The checked condition $M\ge4R+2S$ for $R>0$ is a convenient sufficient
corollary, not the exact threshold.  Moreover, with
$N=2^{M_n}$ and $Q=3^{2R_n}2^{2S_n}$, convexity of the fibre sizes gives at
least
\[
 \frac12\left(\frac{N^2}{Q}-N\right)
\tag{8.2}\label{long1049:eq:collision-count}
\]
unordered equal-signature selector pairs.  A normality estimate can therefore
win by showing that fewer than this many collision pairs land in the two bad
nullspaces.

\begin{problem}[four-jet kernel escape]\label{long1049:prob:escape}
For one literal polynomially primitive family
satisfying~\eqref{long1049:eq:exact-jet-threshold},
prove, preferably by comparing the lower bound~\eqref{long1049:eq:collision-count} with
the bad-pair multiplicities, that
\[
 \mathcal C_n\mathbin{\backslash}
 \bigl(K_n^{\mathrm{poly}}\cup K_n^{\mathrm{rem}}\bigr)
 \ne\varnothing
\]
for all sufficiently large $n$.
\end{problem}

Theorem~\ref{long1049:res:jetkernel} supplies only $\mathcal C_n\ne\varnothing$.
It gives a nonzero selector difference, but it gives neither polynomial-pair
nonvanishing nor remainder nonvanishing.  Problem~\ref{long1049:prob:escape} is finite
algebra and normality; it makes no asymptotic product-formula claim.  It is the
first of the two gates in Problem~\ref{long1049:prob:kernel}: it asks for the two
nonvanishing conclusions, and not for the analytic condition stated there.

\subsection{Determinant families}

An exploratory calculation suggests a contiguous $a_0$-shift determinant at
$r_n=13n+2$, after apparent failure of lower ranks $r\le13n+1$.  These
computations are not Lean theorems.  More importantly, no literal definition
of the matrix $A_{n,r}$ is given here, so writing merely
“$\det A_{n,13n+2}\ne0$” would not yet be a public mathematical question.

\begin{problem}[the first saturated determinant]
Give the entrywise formula for the contiguous $a_0$-shift matrix $A_{n,r}$,
the exact relation $M_n=M(n,r)$ between its size and the number of coefficient
pairs, and the declared width $W_n$.  Then determine whether
\[
 \det A_{n,13n+2}\ne0
\]
for every sufficiently large $n$.  At the first computable scales, report the
rank, determinant, coefficient-pair contents, specialised row contents, four
endpoint jets, and the first nonzero remainder coefficient.  A negative answer
must identify a systematic rank relation or the first exact failing $n$.
\end{problem}

This formulation makes the matrix definition part of the problem statement.  No
undeclared notation is relied on.

\begin{problem}[minimal non-$a_0$ deformation]
If the saturated contiguous family collapses, enlarge it by exactly one
$a_1$-shift while retaining the same source cone and widths.  Does this
one-direction extension increase the polynomial-pair rank and produce a
collision outside
$K_n^{\mathrm{poly}}\cup K_n^{\mathrm{rem}}$?  If every such one-shift
extension collapses, prove that class-wide obstruction before adding a second
new direction.
\end{problem}

This is the smallest specified deformation beyond the contiguous family; it
replaces the unbounded request for a “genuinely independent deformation.”

\subsection{Asymptotic adequacy after kernel escape}

Suppose a good collision has been found.  Let $D_n=3^{R_n}2^{S_n}$ be the
certified local divisor (or replace it by the exact determinant-specific
divisor), let $H_n$ be the coefficient or exterior height after precisely the
normalisation just declared, and let $L_n\ne0$ be the resulting analytic form
or exterior remainder.

\begin{problem}[negative normalised product-formula margin]
Prove the explicit estimate
\[
 \limsup_{n\to\infty}
 \frac{\log H_n+\log|L_n|-R_n\log3-S_n\log2}{n^2}<0.
\tag{8.3}\label{long1049:eq:negative-margin}
\]
Every denominator, row content and final-combination content must already be
included in $H_n$ and $L_n$.
\end{problem}

Nonvanishing alone does not address~\eqref{long1049:eq:negative-margin}.  Conversely,
excellent formal decay is irrelevant if every jet collision lies in a bad
nullspace.  The exploratory adjacent-exterior calculation leaves a positive
normalised exponent of about $110.850\,n^2$; this is not a checked theorem.
Any proposed improvement must remove this explicit deficit; extra divisibility
alone is insufficient.

\subsection{One exact alternative-criterion test}

A separate possible method is Mahler's method.  Its first applicability test has
an exact negative answer, which we record here; an earlier version of this note
left it open.

\begin{proposition}[no finite simultaneous $2/3$-system]\label{long1049:res:nomahler}
Let
\[
 \mathcal L(z)=\sum_{n\ge1}\frac{z^n}{1-z^n}.
\]
There is no finite-dimensional $\Q(z)$-vector space that contains $\mathcal L$
and is stable under both $z\mapsto z^2$ and $z\mapsto z^3$.
\end{proposition}

\begin{proof}
Suppose $V$ were such a space, of dimension $d$.  Stability under $z\mapsto z^2$
places the $d+1$ elements $\mathcal L(z),\mathcal L(z^2),\dots,\mathcal
L(z^{2^{d}})$ in $V$, so they are linearly dependent over $\Q(z)$; clearing
denominators gives polynomials $P_0,\dots,P_d$, not all zero, with
$\sum_{i=0}^{d}P_i(z)\mathcal L(z^{2^i})=0$.  Thus $\mathcal L$ is $2$-Mahler,
and the same argument under $z\mapsto z^3$ makes it $3$-Mahler.  Since $2$ and
$3$ are multiplicatively independent, a theorem of Adamczewski and Bell
\cite[Theorem~1]{adamczewskibell2013} then forces $\mathcal L$ to be a rational
function.  That contradicts Rivin's periodic-coefficient corollary
\cite[Cor.~6.4, p.~9]{rivin2026}, by which $\mathcal L$ is not rational.
\end{proof}

The proposition uses no property of the point $2/3$: the obstruction is
functional and appears before regularity at a particular point is considered.
Bell and Smertnig have since proved the stronger single-base conclusion that
the divisor-function generating series is not $k$-Mahler for any $k\ge2$
\cite{bellsmertnig2026}; thus their result subsumes this simultaneous
$2$-and-$3$ obstruction.  These statements are functional: they do not by
themselves decide the arithmetic nature of the special value
$\mathcal L(2/3)$, nor do they exclude approximation methods outside the
Mahler framework.  The external ingredients are cited rather than proved here.

\subsection{The quantitative height limitation}

The module \texttt{HermitePadeNoGo} defines an explicit rectangular
two-parameter exponent model.  Writing $\sigma=1+\rho+u$, its
denominator-cleared gap has the exact expansion
\[
 -\pi^2\rho^2-\pi^2\rho u-2\pi^2\rho-2\rho^2-10\rho u
 -4\rho-6u^2-8u.
\]
Thus for $\rho\ge0$ and $\sigma\ge1+\rho$ the cleared gap is nonpositive, and
it vanishes exactly when $\rho=0$ and $\sigma=1$
(\lword{Erdos1049/HermitePadeNoGo.lean}{48}{hpClearedGap_expansion}{exact expansion},
\lword{Erdos1049/HermitePadeNoGo.lean}{58}{hpClearedGap_nonpos}{nonpositivity},
\lword{Erdos1049/HermitePadeNoGo.lean}{75}{hpClearedGap_eq_zero_iff}{equality case}).
Equivalently, within that model the displayed threshold never exceeds the
classical one-function margin, with equality only at the classical endpoint
(\lword{Erdos1049/HermitePadeNoGo.lean}{103}{rectangular_hp_threshold_le_classical}{bound},
\lword{Erdos1049/HermitePadeNoGo.lean}{126}{rectangular_hp_threshold_eq_classical_iff}{equality}).
This is an algebraic theorem about the four defined exponent expressions and
only that explicit rectangular model.  It does not construct polynomials,
remainders, integrality, a determinant, or asymptotics, and it is not a
method-universal no-go theorem.
The separate $81/200$ cutoff of the preceding section satisfies
\[
 \frac{81}{200}<\frac{\log2}{\log3}.
\]
Thus “improve the height theorem” has a precise numerical target.

\begin{problem}[optimal admissible height threshold]
Specify a concrete class $\mathcal A$ of primitive, noncollapsed constructions,
excluding scalar row rescaling, proportional permutation-orbit forms, pure
row-content gain, and constructions governed by the rectangular exponent model
above.  For that class define
\[
 \theta_*=\sup_{\alpha\in\mathcal A}
   \frac{C_0(\alpha)}{C_1(\alpha)}.
\]
Prove one of the following: $\theta_*>81/200$ by an explicit construction;
an exact value for $\theta_*$; a converse
$\theta_*<\log2/\log3$; or a class theorem showing that every member reduces to
one of the checked scalar or rectangular no-go mechanisms.
\end{problem}

Reaching $3/2$ requires a threshold strictly beyond
$\log2/\log3\approx0.6309$; exceeding $81/200=0.405$ is insufficient.
The unrestricted question of which rational bases give irrational values
remains open and is not reduced to any one of these problems.

\pdfbookmark[1]{Statements and declarations}{statements-declarations}
\subsection*{Statements and declarations}

\paragraph{Artefact and data availability.} The
\href{\sourceurl}{pinned formal-source revision} contains the Lean sources, the
fixed toolchain, and the library manifest used in the verification.  Proof
authority rests in those pinned sources.  The present text is exposition and
navigation.

\paragraph{Funding and competing interests.} This work received no external
funding.  The author declares no competing interests.

\paragraph{Acknowledgements.} The problem numbering and status follow the
Erd\H{o}s Problems catalogue maintained by Thomas Bloom~\cite{erdosproblems}.

\appendix
\section{Guide to the formal sources}\label{long1049:app:index}

Each linked phrase opens its Lean declaration at the pinned source
revision~\commitshort.  The declarations of this note live in nine modules:
\texttt{RationalBaseLambert}, \texttt{QAperyDiagonalNonEquivalence},
\texttt{RationalPadeArithmetic}, \texttt{ZudilinConeArithmetic},
\texttt{ZudilinHeightRegion}, \texttt{HermitePadeNoGo},
\texttt{BezoutPluckerJets}, \texttt{AdelicHeightBridge}, and
\texttt{RationalBaseContour}.  The first contains
the corridor, cleared-tail recurrence, and elementary $7/2$ certificate; the
second checks the finite $n=0$ diagonal residual; the next five separate
the Pad\'e exponent arithmetic, endpoint arithmetic, logarithmic comparisons,
rectangular exponent model, B\'ezout--Pl\"ucker tail collapse, power
certificates, charge comparisons, four-jet counting and the Hankel assembly;
the last holds the constant, its rational bracket and the degree bookkeeping.
The link
coordinates are validated against that pinned revision, so they remain correct
as later work moves lines in the working tree.

The finite arithmetic behind the rational bracket
$81/200<\theta^{*}<1/2$ of Section~\ref{long1049:sec:regionbracket} and the integer
degree bookkeeping of Section~\ref{long1049:sec:degrees} lives in a further module,
\texttt{RationalBaseContour}, which the library root imports at the pinned
revision~\commitshort.  Its declarations are linked at their line coordinates in
those two sections.

% The linked source manifest is retained in the authored TeX for automated
% validation, and kept outside the printed note.
\iffalse

\begin{tabular}{@{}ll@{}}
informal statement & linked Lean source\\
Van Assche diagonal initial values & \lrefx{Erdos1049/QAperyDiagonalNonEquivalence.lean}{33}{vanAsscheDiagonal_initial_values}\\
exact $n=0$ residual factorisation & \lrefx{Erdos1049/QAperyDiagonalNonEquivalence.lean}{67}{vanAsscheQAperyResidualAtZero_eq}\\
nonzero $n=0$ residual for $p>1$ & \lrefx{Erdos1049/QAperyDiagonalNonEquivalence.lean}{94}{vanAsscheQAperyResidualAtZero_ne_zero}\\
corridor predicate & \lrefx{Erdos1049/RationalBaseLambert.lean}{149}{CoordinatewiseCorridor}\\
corridor bound & \lrefx{Erdos1049/RationalBaseLambert.lean}{157}{coordinatewiseCorridor_implies_pow_lt_linear}\\
exponential comparison & \lrefx{Erdos1049/RationalBaseLambert.lean}{178}{three_mul_lt_two_pow_succ}\\
no corridor at $3/2$ & \lrefx{Erdos1049/RationalBaseLambert.lean}{191}{threeHalves_no_coordinatewiseCorridor}\\
rational-base prefix & \lrefx{Erdos1049/RationalBaseLambert.lean}{204}{rationalBasePrefixQ}\\
prefix step & \lrefx{Erdos1049/RationalBaseLambert.lean}{209}{rationalBasePrefixQ_succ}\\
cleared tail state & \lrefx{Erdos1049/RationalBaseLambert.lean}{217}{rationalBaseClearedTailQ}\\
cleared-tail recurrence & \lrefx{Erdos1049/RationalBaseLambert.lean}{223}{rationalBaseClearedTailQ_succ}\\
forcing magnitude & \lrefx{Erdos1049/RationalBaseLambert.lean}{234}{rationalBaseForcingNat}\\
exponential lower bound & \lrefx{Erdos1049/RationalBaseLambert.lean}{240}{twoPow_le_rationalBaseForcingNat}\\
integer-base collapse & \lrefx{Erdos1049/RationalBaseLambert.lean}{254}{rationalBaseForcingNat_one}\\
integer power comparison & \lrefx{Erdos1049/RationalBaseLambert.lean}{32}{sevenHalves_power_certificate}\\
exact rational boundary & \lrefx{Erdos1049/RationalBaseLambert.lean}{36}{sevenHalves_rational_margin}\\
logarithmic ratio bound & \lrefx{Erdos1049/RationalBaseLambert.lean}{48}{sevenHalves_log_ratio_lt}\\
$\pi$ reciprocal-square bound & \lrefx{Erdos1049/RationalBaseLambert.lean}{64}{one_div_pi_sq_lt_one_nine}\\
strict source margin & \lrefx{Erdos1049/RationalBaseLambert.lean}{73}{sevenHalves_bundschuhVaananen_margin}\\
complete $7/2$ height condition & \lrefx{Erdos1049/RationalBaseLambert.lean}{119}{sevenHalves_archimedean_height_condition}\\
doubled common exponent & \lrefx{Erdos1049/RationalPadeArithmetic.lean}{19}{rationalPadeDenExpTwice}\\
doubled summand exponent & \lrefx{Erdos1049/RationalPadeArithmetic.lean}{24}{rationalPadePSummandDenExpTwice}\\
summand exponent bound & \lrefx{Erdos1049/RationalPadeArithmetic.lean}{30}{rationalPadePSummandDenExpTwice_le}\\
doubled maximal exponent & \lrefx{Erdos1049/RationalPadeArithmetic.lean}{46}{rationalPadeQMaxDenExpTwice}\\
exact gap identity & \lrefx{Erdos1049/RationalPadeArithmetic.lean}{52}{rationalPadeQMaxDenExpTwice_gap}\\
maximal exponent bound & \lrefx{Erdos1049/RationalPadeArithmetic.lean}{60}{rationalPadeQMaxDenExpTwice_le}\\
Pad\'e error scaling & \lrefx{Erdos1049/RationalPadeArithmetic.lean}{84}{rationalPadeError_mul}\\
row-content determinant factorisation & \lrefx{Erdos1049/RationalPadeArithmetic.lean}{94}{rationalPadeExteriorDet_mul_contents}\\
absolute determinant scaling & \lrefx{Erdos1049/RationalPadeArithmetic.lean}{104}{natAbs_rationalPadeExteriorDet_mul_contents}\\
content-product divisibility & \lrefx{Erdos1049/RationalPadeArithmetic.lean}{116}{contentProduct_dvd_rationalPadeExteriorDet}\\
exterior determinant identity & \lrefx{Erdos1049/RationalPadeArithmetic.lean}{124}{rationalPadeExteriorDet_cast_eq}\\
homogeneous endpoint evaluation & \lrefx{Erdos1049/ZudilinConeArithmetic.lean}{98}{homEvalThreeTwo}\\
four-jet target cardinality & \lrefx{Erdos1049/ZudilinConeArithmetic.lean}{131}{fourJetSignature_card}\\
four-jet pigeonhole kernel & \lrefx{Erdos1049/ZudilinConeArithmetic.lean}{142}{exists_distinct_binary_selectors_same_fourJet}\\
rank--depth collision threshold & \lrefx{Erdos1049/ZudilinConeArithmetic.lean}{161}{exists_distinct_binary_selectors_same_fourJet_of_rank}\\
bottom-jet divisibility & \lrefx{Erdos1049/ZudilinConeArithmetic.lean}{191}{bottomJet3_eq_zero_iff_dvd}\\
top-jet divisibility & \lrefx{Erdos1049/ZudilinConeArithmetic.lean}{197}{topJet2_eq_zero_iff_dvd}\\
bottom endpoint modulo three & \lrefx{Erdos1049/ZudilinConeArithmetic.lean}{203}{homEvalThreeTwo_mod_three}\\
generic homogeneous evaluation & \lrefx{Erdos1049/ZudilinConeArithmetic.lean}{221}{homEval}\\
left endpoint congruence & \lrefx{Erdos1049/ZudilinConeArithmetic.lean}{226}{homEval_mod_left}\\
right endpoint congruence & \lrefx{Erdos1049/ZudilinConeArithmetic.lean}{238}{homEval_mod_right}\\
left endpoint coprimality & \lrefx{Erdos1049/ZudilinConeArithmetic.lean}{254}{homEval_isCoprime_left_of_const_unit}\\
right endpoint coprimality & \lrefx{Erdos1049/ZudilinConeArithmetic.lean}{270}{homEval_isCoprime_right_of_top_unit}\\
two-endpoint coprimality & \lrefx{Erdos1049/ZudilinConeArithmetic.lean}{286}{homEval_isCoprime_mul_of_endpoint_units}\\
cyclotomic endpoint coprimality (Proposition~3.6) & \lrefx{Erdos1049/ZudilinConeArithmetic.lean}{302}{cyclotomicHomEval_isCoprime_mul}\\
top endpoint modulo two & \lrefx{Erdos1049/ZudilinConeArithmetic.lean}{320}{homEvalThreeTwo_mod_two}\\
unit constant endpoint obstruction & \lrefx{Erdos1049/ZudilinConeArithmetic.lean}{338}{three_not_dvd_homEvalThreeTwo_of_const_unit}\\
unit top endpoint obstruction & \lrefx{Erdos1049/ZudilinConeArithmetic.lean}{356}{two_not_dvd_homEvalThreeTwo_of_top_unit}\\
common-multiplier exclusion & \lrefx{Erdos1049/ZudilinConeArithmetic.lean}{398}{commonMultiplier_not_two_not_three_of_endpoint_units}\\
irrationality-exponent inequality & \lrefx{Erdos1049/ZudilinConeArithmetic.lean}{416}{twice_le_of_irrationalityExponent_bounds}\\
three-halves scalar margin & \lrefx{Erdos1049/ZudilinConeArithmetic.lean}{428}{three_two_scalar_margin_neg}\\
$31/4$ logarithmic-region membership & \lrefx{Erdos1049/ZudilinHeightRegion.lean}{45}{thirtyoneFour_mem_zudilinHeightRegion}\\
positive-power membership & \lrefx{Erdos1049/ZudilinHeightRegion.lean}{59}{thirtyoneFour_power_mem_zudilinHeightRegion}\\
$31/4$ outside the earlier region & \lrefx{Erdos1049/ZudilinHeightRegion.lean}{91}{thirtyoneFour_outside_bundschuhVaananenHeightRegion}\\
$81/200$ below the $3/2$ log ratio & \lrefx{Erdos1049/ZudilinHeightRegion.lean}{105}{eightyOneTwoHundredths_lt_threeHalves_log_ratio}\\
$3/2$ outside the $81/200$ region & \lrefx{Erdos1049/ZudilinHeightRegion.lean}{122}{threeHalves_outside_zudilinHeightRegion}\\
rectangular exponent-model bound & \lrefx{Erdos1049/HermitePadeNoGo.lean}{103}{rectangular_hp_threshold_le_classical}\\
equality case for the exponent model & \lrefx{Erdos1049/HermitePadeNoGo.lean}{126}{rectangular_hp_threshold_eq_classical_iff}\\
\end{tabular}
\fi

\paragraph{Checked declarations behind the four-jet collision family.}
For a modular tail whose every row is unimodular and whose adjacent minors
vanish, all minors vanish,
\lword{Erdos1049/BezoutPluckerJets.lean}{191}{adjacent_det_zero_forces_all_det_zero_of_isCoprime}{adjacent-minor
propagation}, the sharper $N<2^k$ collision threshold applies,
\lword{Erdos1049/BezoutPluckerJets.lean}{280}{zmod_binary_tail_collision_of_adjacent_det_zero_of_isCoprime}{binary
tail collision}, and at modulus $2^S3^R$ with $R>0$ the depth $k\ge S+2R$
suffices,
\lword{Erdos1049/BezoutPluckerJets.lean}{299}{zmod_binary_tail_collision_of_two_three_depth_of_isCoprime}{two-three
depth collision}.  Neither coordinate must be a unit; $(2,3)$ modulo six is an
example.
